\documentclass[
  ngerman,
  ruledheaders=section,%Ebene bis zu der die Überschriften mit Linien abgetrennt werden, vgl. DEMO-TUDaPub
  class=report,% Basisdokumentenklasse. Wählt die Korrespondierende KOMA-Script Klasse
  %thesis={type=bachelor},% Dokumententyp Thesis (projektseminar, proseminar, bachelor, master)
  accentcolor=2c, % Akzentfarbe FG rtm
  custommargins=false,
  marginpar=false,
  BCOR=10mm, % Bindekorrektur (etwas größere mittlere Ränder)
  parskip=half-, % Absatzkennzeichnung durch Abstand vgl. KOMA-Sript
  fontsize=11pt, % Basisschriftgröße laut Corporate Design ist mit 9pt häufig zu klein
  twoside,
  numbers=noendperiod,
  fleqn,
  captions=oneline, % Einzeilige Captions zentrieren
  captions=tableheading,
  IMRAD=false,
  pdfa=false,
]{tudapub}


% Der folgende Block ist nur bei pdfTeX auf Versionen vor April 2018 notwendig
\usepackage{iftex}
\ifPDFTeX
\usepackage[utf8]{inputenc}%kompatibilität mit TeX Versionen vor April 2018
\fi

%%%%%%%%%%%%%%%%%%%
%Sprachanpassung & Verbesserte Trennregeln
%%%%%%%%%%%%%%%%%%%
\usepackage[english, main=ngerman]{babel}
\usepackage[autostyle]{csquotes}% Anführungszeichen vereinfacht
\usepackage{microtype}


%%%%%%%%%%%%%%%%%%%
%Literaturverzeichnis
%%%%%%%%%%%%%%%%%%%
% ACHTUNG: bei Problemen ggf. den ieee Style nachinstallieren oder einen anderen nutzen
\catcode30=12  % Workaround for current biblatex problem (https://tex.stackexchange.com/questions/564990/error-after-miktex-reinstall-text-line-contains-an-invalid-character)
\usepackage[backend=biber,style=ieee,bibencoding=utf8]{biblatex}
\addbibresource{literature.bib}


\input{common/includes.tex}

\input{common/include_and_setup_lstlistings.tex}
\input{common/include_and_setup_pgf.tex}

\input{common/commonmacros.tex}


\newcommand*{\Miktex}{\textsc{MiKTeX}\xspace}
\newcommand*{\texlive}{\textsc{texlive}\xspace}
\newcommand{\Texstudio}{\textsc{TeXStudio}\xspace}
\newcommand{\Sumatra}{\textsc{Sumatra}\xspace}
\newcommand*{\zitat}[1]{\glqq{}#1\grqq{}}



% \Metadata{  % Metadaten im erzeugten PDF/A
%   title=Hinweise zum Abfassen einer (Abschluss-)Arbeit am IAT,
%   author=Rudi Ratlos,
%   keywords={TU Darmstadt \sep Corporate Design \sep LaTeX},
% }

\title{Hinweise zum Verfassen einer (Abschluss-)Arbeit am IAT}
%\author{Rudi Ratlos}
\date{Stand: 21.10.2021}
%\reviewer{Prof. Dr.-Ing. Ulrich Konigorski \and Wi Mi M.Sc.}

%\department{etit}

\addTitleBox{\includegraphics[width=\linewidth]{CCPS_logo}}

%\date{\today}
%\submissiondate{1. April 2020}

\lowertitleback{
  Technische Universität Darmstadt\\
  Institut für Automatisierungstechnik und Mechatronik\\
  Fachgebiet Regelungstechnik und Mechatronik\\
  Prof. Dr.-Ing. Ulrich Konigorski\\
}

\begin{document}

\frontmatter
\maketitle

\cleardoublepage
\tableofcontents

\cleardoublepage

\mainmatter

\chapter{Einführung}
\label{cha:Intro}

Beim Anfertigen einer Abschlussarbeit steht man als Studierender meist zum ersten Mal vor dem Problem, einen längeren wissenschaftlichen Text mit Bildern, Gleichungen und Referenzen schreiben zu müssen.
Dafür bietet sich das Textsatz-System \LaTeX\ an, zu dessen Vorteilen die weitgehende Trennung von Inhalt und Layout gehören.

Die Arbeiten sollen im Corporate Design der TU Darmstadt abgefasst werden.
Hierzu stellt die TU Darmstadt das \LaTeX-Paket \verb|tuda-ci| zur Verfügung, welches auch für den vorliegenden Text verwendet wurde.


Allgemeines zu den Grundlagen von \LaTeX{} findet man \zB in \cite{Kopka} oder \cite{Schmidt}.
Allgemeines über wissenschaftliche Arbeiten findet man in \cite{Friedrich}.
Als Einstieg in die Typografie ist \cite{Willberg} sehr zu empfehlen.

Die Arbeit kann in Deutsch oder wahlweise auch in Englisch verfasst werden.

In der vorliegenden Anleitung werden in \charef{cha:Hinweise} zuerst allgemeine Hinweise und Tipps zur Durchführung einer wissenschaftlichen Arbeit gegeben.
Im Anschluß daran geht \charef{cha:Hinweise_Latex} kurz auf die Verwendung von \LaTeX{} zum Erstellen der schriftlichen Arbeit ein.
In \charef{cha:Hinweise_Allgemein} wird dann wieder allgemein auf die Abfassung schriftlicher Arbeiten eingegangen.
Dazu wird zu jedem Punkt nach den allgemeinen Hinweisen auch die Umsetzung in \LaTeX{} angesprochen.

\charef{cha:Beurteilung} umreißt die Kriterien, nach denen eine Arbeit an unserem Fachgebiet bewertet wird.

Im Anhang findet sich eine Checkliste die vor Abgabe der Arbeit unbedingt abgearbeitet werden sollte, um zu prüfen, ob alle Vorgaben und Hinweise beachtet wurden.

Dieses Dokument ist \emph{keine} Einführung in \LaTeX{}.
Es werden zwar einige Befehle und Umgebungen namentlich genannt, jedoch sind diese Aufzählungen weder vollständig noch werden die Befehle hier genauer erläutert.
Hierzu sollten sich im Internet schnell alle Informationen finden lassen.
Erfahrungsgemäß ist es auch nicht nötig, spezielle Tutorials oder Kurse zu \LaTeX\ durchzuführen.
Ausgehend von dem zur Verfügung gestellten Minimalbeispiel lässt sich das notwendige Wissen in der Regel gut iterativ, während der Bearbeitung der Arbeit, erlangen.
Empfehlenswert ist es dazu, möglichst früh mit \LaTeX{} zu beginnen, und beispielsweise schon Stichworte in der Einarbeitungsphase der Arbeit zu "`texen"'.

\chapter{Tipps für die Durchführung einer wissenschaftlichen Arbeit}
\label{cha:Hinweise}

Die Bearbeitung einer wissenschaftlichen (Abschluss-)Arbeit gliedert sich in vier Phasen, die im Folgenden anhand von Stichworten umrissen werden.

\section{Phase 1: Einarbeitung und Literaturrecherche}
\label{sec:Einarbeitung}

\begin{itemize}
  \item dient zur Verdeutlichung der gegebenen Problemstellung
  \item führt den Bearbeiter auf den \glqq{}State-of-the-Art\grqq{} hin $\rightarrow$ Literatur
  \item zeigt Möglichkeiten zur Problemlösung auf, die dann in Phase 2 genutzt werden können
  \item dient der Ideenfindung
  \item Erstellen eines Zeitplanes
  \item regelmäßige Treffen mit dem Betreuer
\end{itemize}



\subsection*{Literaturrecherche und -management}
\label{sec:Literaturrecherche}
Zur Sammlung und Verwaltung von Literatur dienen sog. Literaturdatenbanken.
Diese Programme ermöglichen oft die direkte Verwendung als Literaturverzeichnis im späteren Text.

Das Suchen und Lesen der aktuellsten Literatur ist wichtig für eine gute Übersicht über die Problemstellung.
Über die ULB ist es möglich viele Online-Quellen%
\footnote{Für einige ist es nötig, sich im TUD-Netz zu befinden, ggf. per VPN.
Hierzu finden sich weitere Infos auf der TUD-Website.}
zu nutzen, um Texte auch als pdf-Dokumente erhalten zu können.
Die für unseren Bereich wichtigsten Seiten sind:
\begin{itemize}
  \item \href{https://scholar.google.de/}{Google Scholar}
  \item \href{http://rzblx1.uni-regensburg.de/ezeit/}{Elektronische Zeitschriftenbibliothek} (DB aller Zeitschriften und Berechtigungen)
  \item \href{http://ieeexplore.ieee.org}{IEEEXplore} (Zugang zu allen Medien der IEEE)
  \item \href{http://www.sciencedirect.com}{Sciencedirect} (Zugang zu Zeitschriften aus dem ELSEVIER Verlag)
  \item \href{http://citeseerx.ist.psu.edu}{CiteSeer} (freie DB mit Zugang zu vielen Volltexten)
  \item \href{http://www.ifac-control.org/publications/ifac-papersonline.net}{IFAC-PapersOnline} (Archiv mit Artikeln aller IFAC Konferenzen)
  \item \href{http://link.springer.com/}{Springer E-Books} (Archiv des Springer Verlags mit vielen eBooks und Zeitschriften)
\end{itemize}
auch mal auf den Webseiten der Autoren schauen und natürlich \glqq{}googlen\grqq{}.


\section{Phase 2: Bearbeitung des Problems}
\label{sec:Bearbeitung}

\begin{itemize}
  \item dient zur Lösung der Problemstellung
  \item kreativste und anstrengendste Phase der Arbeit
  \item Verwendung professioneller Hilfsmittel (Programme wie Matlab oder Mathematica etc.)
  \item Treffen wenn Bedarf besteht, keine Regelmäßigkeit mehr
  \item sowohl die Implementierung als auch die Ausarbeitung/Präsentation sollten spätestens ab dieser Phase mittels eines Versionskontrollsystems wie \bspw \textsc{Git} versioniert werden und regelmäßig zur Datensicherung auf einen externen Speicher oder einen Cloud-Speicher (\zB Hessenbox) übertragen werden, damit ältere (möglicherweise noch funktionsfähige) Stände wiederhergestellt werden können und keine Daten bei Zerstörung oder Verlust des Arbeitsrechners verloren gehen\\
    \textbf{Beachten Sie:}
    Sie als Bearbeiter der Arbeit sind für die Datensicherung selber verantwortlich.
    Dies gilt auch, wenn Sie an einem Rechner des Instituts (Rechnerpool oder Versuchstand) arbeiten!
  \item Eine gründliche Dokumentation von Teilergebnissen (Herleitungen, Rechenwege, verwendete Nomenklatur, \zB bereits in Latex) erleichtert später das Verfassen der Arbeit.
\end{itemize}


\section{Phase 3: Ergebnisse zusammenstellen}
\label{sec:Ergebnisse}

\begin{itemize}
  \item dient der Reorganisation der Arbeit
  \item sämtliche Ergebnisse werden festgehalten
  \item Strukturierung und Gliederung der Ergebnisse, so dass die nächste Phase (Schreiben) gut durchgeführt werden kann
  \item wenige, längere Treffen zur Ergebnisbesprechung mit dem Betreuer
\end{itemize}


\subsection*{Ergebnissicherung}
\label{sec:Ergebnissicherung}
\begin{itemize}
  \item alles zusammentragen was erreicht wurde $\rightarrow$ guter Überblick notwendig (s.\,o.)
  \item auch Programmcode ist ein Ergebnis $\rightarrow$ verständlich kommentieren (Englisch)
  \item auf  Wiederverwendbarkeit von Grafiken achten (Linienstärke, Farbe, Beschriftung, \ldots)
  \item nur noch kleine Änderungen durchführen (z.B. Parametereinstellungen)
\end{itemize}
\emph{Ergebnisse sollten für sich sprechen und für jeden verständlich sein (ohne Erklärung)}


\section{Phase 4: Dokumentation und Präsentation}
\label{sec:Dokumentation}

\begin{itemize}
  \item Verfassen der Arbeit und Erstellen der Präsentation
  \item letzte Verfeinerungen an den Ergebnissen (falls notwendig)
  \item schwierigste Phase der Arbeit
\end{itemize}


\subsection*{Wissenschaftliches Schreiben}
\label{sec:Wissenschaftliches Schreiben}
\begin{itemize}
  \item nicht am Layout der Arbeit aufhalten $\rightarrow$ Vorlage verwenden
  \item sehr aufwendiger Prozess von Schreiben -- Verbessern -- neu Schreiben -- \ldots
  \item Aufwand darf nicht unterschätzt werden (Richtwert: 1-2 Seiten/Tag)
  \item einfach erst einmal aufschreiben --  korrigiert wird dann später
  \item Struktur einer wissenschaftlichen Arbeit ist vorgegeben
    \begin{itemize}
      \item \textbf{Titelseite} mit Art der Arbeit, Titel, Namen des Autors sowie Abgabedatum.
      \item \textbf{Aufgabenstellung} wird vom Betreuer der Arbeit zur Verfügung
        gestellt.
      \item \textbf{Erklärung} zur Selbständigkeit.
        Der Text ist vorgegeben und wird bei Verwendung der Vorlage automatisch erzeugt.
      \item \textbf{Kurzfassung} der Arbeit.
        Der Umfang soll so bemessen sein, dass die englische Version (\textbf{Abstract}) auf die gleiche Seite passt.
      \item \textbf{Inhaltsverzeichnis} wird in \LaTeX{} durch \verb|\tableofcontents| automatisch erzeugt.
      \item \textbf{Symbole und Abkürzungen}.
        Dieses Verzeichnis erstellt man am Besten von Hand.
        Die Einteilung in lateinische und griechische Symbole und Formelzeichen kann nach Bedarf geändert werden (zum Beispiel nach Kapiteln oder Konzepten) oder ganz weggelassen werden.
      \item \textbf{Hauptteil der Arbeit}, in einzelne Kapitel und Abschnitte unterteilen.
      \item \textbf{Anhang}.
        Hier können Abschnitte stehen, die beim Lesen der Arbeit stören würden, \zB\ Programmcode, technische Daten oder lange mathematische Beweise.
      \item \textbf{Literaturverzeichnis} sollte automatisch generiert (in \LaTeX{} \zB mit biblatex/biber). Formatierung und Inhalt prüfen und ggf. korrigieren!
      \item \label{itm:Ordner} \textbf{Zusätzliches Material} wie \zB\ der vollständige Programmcode eines Software-Projekts gehört nicht in die Arbeit, sondern kann in einem separaten Ordner abgelegt werden.
    \end{itemize}
  \item Wir bieten Ihnen an, am Ende einmal Ihre Arbeit korrektur zu lesen.
    Darüberhinaus können Sie sich mit konkreten Fragen zu einzelnen Aspekten natürlich auch immer an Ihren Betreuer wenden.
  \item Orientieren Sie sich beim Schreiben an "`POEM"':\\
    $P_\text{rägnanz}\; O_\text{rdnung}\; E_\text{infachheit}\; M_\text{otivation}$ (einfache, kurze, strukturierte und anregende Sätze)
\end{itemize}


\paragraph{Inhalt und Umfang der Arbeit}
Ihre Arbeit soll am Ende \emph{vollständig} sein, und dies hat zunächst nichts mit der Länge in Seiten zu tun, sondern mit inhaltlichen Fragen:
\begin{itemize}
  \item Wird die Problemstellung und Voraussetzungen der Arbeit klar?
  \item Kann man Ihrem Lösungsansatz folgen?
  \item Wird das Ergebnis ausreichend diskutiert und eingeordnet, \zB durch Vergleiche mit alternativen Ansätzen oder bestehenden Methoden?
  \item Wird bei negativen Ergebnissen eine überzeugende Analyse durchgeführt, warum ein gewählter Ansatz nicht funktioniert hat?
\end{itemize}
Je nach Art der Arbeit umfasst eine vollständige Angabe noch weitere Aspekte, die nicht unbedingt Teil der eigentlichen Ausarbeitung sind:
\begin{itemize}
  \item Sind (bei konstruktiven Arbeiten) die Konstruktionsunterlagen vollständig?
  \item Ist Programmcode, der weiterverwendet werden soll, entsprechend ausführlich kommentiert?
  \item Gibt es zu dem entwickelten Versuchsaufbau eine verständliche Anleitung?
\end{itemize}

Es soll an dieser Stelle bewusst kein erwarteter Umfang der Arbeit in Seiten angegeben werden, da dieser je nach der Art der Arbeit (Bachelorarbeit, Masterarbeit, \ldots) und Ausprägung (theoretisch, simulativ, experimentell, konstruktiv) der Arbeit zu stark schwankt.
Bei Fragen und Unsicherheiten zum Umfang der Arbeit wird Ihnen daher Ihr Betreuer weiterhelfen.





\subsection*{Verteidigung und Präsentation}
\label{sec:Verteidigung}
\begin{itemize}
  \item Dauer der Präsentation (ohne Fragerunde): 20 Minuten (bei Bachelor- und Masterarbeiten)
    \begin{itemize}
      \item nach 25 Minuten wird abgebrochen
      \item Grundsätzlich lieber etwas zu lang (aber auf jeden Fall unter 25 Minuten) als zu kurz.
      \item Bei Proseminaren (10 Minuten) und Projektseminaren (je nach Anzahl der Teilnehmer) abweichende Länge
    \end{itemize}
  \item Inhalt der Arbeit auf ca. 10 wesentliche Punkte reduzieren
  \item je Punkt eine Folie, maximal zwei
  \item je Folie 1-2 min Gesprächszeit
  \item Struktur der Präsentation:
  \begin{itemize}
    \item Motivation / Einleitung (Aufgabenstellung)
    \item Grundlagen
    \item Lösungsweg
    \item Ergebnisse
    \item Zusammenfassung (und Ausblick)
  \end{itemize}
  \item Beachten Sie: Was sich durch ein Bild (anstelle von Text) darstellen lässt, lässt sich durch den Zuhörer schneller erfassen und wertet die Präsentation optisch auf.
  \item insbesondere komplexe Sachverhalte durch Abbildungen verdeutlichen
  \item nur Stichpunkte schreiben, keine ganzen Sätze
  \item klare, einfach nachvollziehbare Notation (\zB bei Variablen) verwenden
  \item mehrmaliges Üben des Vortrags (\zB vor dem Spiegel, vor der Familie, \ldots)
  \item Probevortrag vor Betreuer halten
\end{itemize}
\emph{nichts weglassen was auf einer Folie steht, gerne zusätzliche Dinge erwähnen}


\section{Allgemeine Punkte}

\subsection*{Organisatorisches}
\label{sec:Organisatorisches}

\begin{itemize}
  \item jeder Studierende sollte an mindestens einem Regelungstechnischen Seminar teilnehmen
  \item Abgabetermin ist ein fixer Termin $\rightarrow$ Prüfungsleistung (Durchfallen möglich)
  \item jeder Studierende erhält auf Wunsch einen eigenen Account für unsere Rechner (bei Betreuer melden)
  \item zur Erstellung der Ausarbeitung und der Präsentation wird die Verwendung des Textsatzprogramms \LaTeX{} empfohlen, das Template für die Ausarbeitung erhalten Sie von Ihrem Betreuer
\end{itemize}




\subsection*{Generelles zu Treffen mit dem Betreuer}
\label{sec:TreffenBetreuer}

Die persönlichen Treffen mit dem Betreuer sollten gut vorbereitet werden (\zB kurze Präsentation, Fragenkatalog).
Die Ergebnisse der Treffen und die Anmerkungen des Betreuers sollten Sie in Form von Notizen festhalten.


Hier ist zu bemerken, dass ein häufigeres Treffen mit dem Betreuer alleine keine mangelnde Selbständigkeit impliziert.
Solange diese Treffen derart ablaufen, dass Sie Ihren Arbeitsfortschritt seit dem letzten Treffen auf \zB ein paar kurzen Folien vorstellen (die nicht übermäßig "`schön"' gestaltet sein sollen), Probleme benennen, die dabei möglicherweise aufgetreten sind und dann noch erläutern, wie Sie mit diesen umgehen wollen, wird es auch bei einem wöchentlichen Treffen hieraus keine Abzüge in der Selbständigkeit geben.
Wenn Sie dagegen der Auf"|fassung sind, Ihre Selbständigkeit dadurch unter Beweis zu stellen, dass Sie monatelang nichts von sich hören lassen, ist dies hinsichtlich der Kommunikation mit dem Betreuer nicht ideal, und vor allem tragen Sie dann das volle Risiko, wenn Sie die Aufgabe falsch verstanden haben und damit in die falsche Richtung arbeiten!




\subsection*{Allgemeines zu Hilfsmitteln bei der Erstellung und Bearbeitung}
\label{sec:Allgemeines}

\begin{itemize}
  \item es gibt viele Bücher die bei der Erstellung von wissenschaftlichen Arbeiten helfen
  \item viele verwendete Programme bieten ausführliche Hilfen an (\zB \Matlab, Mathematica)
  \item bei Problemen immer zuerst in der Hilfe schauen, dann "`googlen"', dann Betreuer fragen
\end{itemize}

\chapter{Hinweise zur Verwendung von \LaTeX}
\label{cha:Hinweise_Latex}

Dieses Kapitel gibt allgemeine Hinweise zur Erstellung einer wissenschaftlichen Arbeit mit dem Textsatzprogramm \LaTeX.


\section{Vorlagen}

Die Arbeit soll auf dem Paket \verb|tuda-ci|, welches von der TU Darmstadt zur Verfügung gestellt wird, aufbauen.
Dies erleichtert es, sich an die Layout-Vorgaben sowie an die formalen Vorgaben (Selbständigkeitserklärung in der aktuellen Fassung) zu halten.

Basierend auf diesem Paket stellt das Fachgebiet rtm verschiedene Vorlagen zur Verfügung, um den Einstieg zu erleichtern:
\begin{itemize}
  \item \verb|tuda-ci_vorlage|\\
    Minimalbeispiel und Startpunkt für die eigene Arbeit.
    Bindet die notwendige Dokumentenklasse aus \verb|tuda-ci| ein, stellt das Fachgebietslogo zur Verfügung und gibt die minimale Grundstruktur der Arbeit vor.\\
    Im Quelltext ist die Anpassung auf verschiedene Varianten einer studentischen Arbeit (Proseminar, externe Arbeit in Unternehmen) erläutert.
  \item \verb|tuda-ci_thesisinfo|\\
    Quelltext zu diesem Dokument.
    Basiert ebenfalls auf \verb|tuda-ci| und zeigt exemplarisch, wie ein etwas größeres Dokument mit mehreren Dateien in einer Ordnerstruktur organisiert werden könnte.
    Man kann sich hieran orientieren, es ist aber nicht als Grundlage der eigenen Arbeit vorgesehen.
  \item \verb|tuda-ci_useful-packages|\\
    Hier sind (minimale) Konfigurationen und die Verwendung ein paar wahrscheinlich nützlicher Pakete (\zB \verb|tikz|) demonstriert.
\end{itemize}



\section{\TeX-System}

Ziel ist es, aus dem in \LaTeX{} geschriebenen "`Quelltext"' ein PDF zu erzeugen. Hierzu wird ein \TeX-Programm, quasi der "`Compiler"', benötigt.
\begin{itemize}
  \item pdflatex\\
    Das meist am Institut verwendete Programm.
    Erzeugt direkt eine \verb|pdf|-Datei und kann dabei direkt \verb|png|- und \verb|jpg|-Bilder verarbeiten.
  \item lualatex\\
    Etwas neueres Programm.
    Erzeugt auch direkt eine \verb|pdf|-Datei.
    Für große Dateien etwas langsamer.
  \item xelatex\\
    Weitere Alternative.
    Erzeugt auch direkt eine \verb|pdf|-Datei.
    (Kann direkt auf die Systemschriftarten zugreifen.)
  \item latex\\
    Das "`ursprüngliche"' Latex-Programm.
    Erzeugt eine Ausgabe im \verb|dvi|-Format, welches dann erst in eine \verb|ps|-Datei und diese dann über Ghostscript in eine \verb|pdf|-Datei umgewandelt werden muss.
    Sollte nicht mehr verwendet werden.
\end{itemize}
In der Anwendung sollten die ersten beiden Möglichkeiten kaum einen Unterschied machen.

Neben dem eigentlichen \LaTeX-Programm werden im Allgemeinen noch weitere Programme benötigt, \zB \verb|biber| zur Erzeugung des Literaturverzeichnisses.

Ein \LaTeX-Programm erzeugt bei einem Durchlauf (Aufruf) eine Menge an Hilfsdateien, die dann weitere Programme oder auch \LaTeX{} selber (in weiteren Durchläufen) verwenden, um beispielsweise Inhaltsverzeichnisse, Literaturverzeichnisse und Referenzierungen zu generieren.
\Dah typischerweise wäre eine Aufrufreihenfolge
\begin{verbatim}
  pdflatex -> biber -> pdflatex -> pdflatex
\end{verbatim}
Um dies zu vereinfachen, kann ein geeigneter Editor und/oder das Programm \verb|latexmk| verwendet werden.

Neben diesen auszuführenden Programmen existieren noch eine sehr große Menge an frei verfügbaren Paketen, die die verschiedensten Befehle zur Verfügung stellen.
(Erzeugung von Grafiken, besonderer Formelsatz, besondere Tabellen, aufrechte griechische Kleinbuchstaben, \ldots)

Zusammen bilden diese Bestandteile das \TeX-System.


\section{\LaTeX-Distribution}
\label{sec:Distribution}

Damit nicht alle oben genannten Komponenten selber und einzeln geladen und installiert werden müssen, sollte auf sogenannte \LaTeX-Distributionen zurückgegriffen werden.
Diese kümmern sich idealerweise um alles, und -- je nach Konfiguration -- laden diese auch noch fehlende Pakete automatisch herunter, wenn diese zum ersten Mal verwendet werden.

Welche Distribution verwendet wird, hängt vom persönlichen Geschmack ab.
Es besteht grundsätzlich die Wahl zwischen \Miktex{}, das vor allem unter Windows weit verbreitet ist und \texlive{}, das eher unter Linux und Mac Anwendung findet.
Daneben gibt es mittlerweile auch Cloud-Lösungen (overleaf), die sich besonders für Gruppenarbeiten eignen, da mehrere Personen gleichzeitig am gleichen Dokument arbeiten können.

Zuerst ist die Distribution herunterzuladen und zu installieren, wobei im Allgemeinen die Standardeinstellungen ausreichend sind.

Das sollte dann eigentlich schon ausreichen, um aus den zur Verfügung gestellten (Minimal-)Beispielen und einem geeigneten Editor (siehe unten) die \verb|pdf|-Dateien zu erzeugen.


\section{Editor}
Für \LaTeX{} gibt es eine Vielzahl von freien und kommerziellen Texteditoren.
Das \Texstudio ist ein freier open-source Editor, der sich bei uns am Institut bewährt hat.
Alternativ kann beispielsweise auch Visual Studio Code mit der "`LaTeX Workshop"'-Erweiterung verwendet werden.

Viele Editoren verfügen über eingebettete PDF-Reader \bzw arbeiten gut \zB mit SumatraPDF zusammen.
\Dah das erzeugte PDF kann im Reader geöffnet bleiben%
\footnote{Dies ist auch der Hauptgrund, weshalb von der Verwendung des Adobe-Acrobat-Readers abgeraten wird.}%
, und man kann aus dem Quellcode an die entsprechende Stelle im PDF und umgekehrt springen.




\section{Verwenden von \LaTeX-Paketen}

Durch das Einbinden von Zusatzpaketen kann \LaTeX\ angepasst und erweitert werden.
Die Pakete werden mit dem Befehl \verb|\usepackage{...}| eingebunden, \ggf können auch noch Optionen in \verb|[...]| angegeben werden.


\section{Definition eigener Befehle}

Die Möglichkeit eigene Befehle in \LaTeX{} zu definieren und zu verwenden erleichtert das Erstellen einer wissenschaftlichen Arbeit deutlich.
So dient ein eigener Befehl oft dazu, häufig verwendete Befehlsfolgen kürzer und schneller schreiben zu können.
Von zentraler Bedeutung ist außerdem, dass man diesen Befehl einfach ändern kann.
Hat man \zB alle Matrizen mit einem eigenen Befehl versehen, der diese fett formatiert, so lässt sich dies auch schnell für alle Matrizen wieder ändern.
Entscheidet man sich Matrizen mit einem Unterstrich zu kennzeichnen, so ist lediglich die Anpassung des entsprechenden Befehls notwendig.

Dies lässt sich auch auf Variablennamen übertragen.
Definiert man \zB für $\tilde{\hat{\mat{x}}}_\mathrm{b2}$ einen neuen Befehl \verb|\xb2|, so verkürzt sich zum einen der Schreibaufwand.
Zum anderen lässt sich selbst in der Endphase der Arbeit die Variable umbenennen, \zB in $\mat{z}_2$, indem lediglich der Befehl verändert wird.

\chapter{Allgemeine Hinweise und Informationen zum Erstellen einer schriftlichen Arbeit}
\label{cha:Hinweise_Allgemein}

Dieses Kapitel gibt allgemeine Hinweise zur Erstellung einer wissenschaftlichen Arbeit und ist damit auch für Studierende relevant, die sich gegen eine Erstellung der Arbeit in \LaTeX{} entschieden haben.
Zusätzlich wird aber auch darauf eingegangen, wie die genannten Punkte mit dem Textsatzprogramm \LaTeX{} umgesetzt werden können.



\section{Quellenarbeit (Zitate und korrekte Zitierweise)}
\label{sec:Latex-Zitieren}



Zitate sind wörtliche oder \emph{auch sinngemäße} Wiedergaben von Gedanken, Ideen oder Meinungen anderer Autoren.
Werden Ideen oder Inhalte aus Quellen wörtlich oder sinngemäß in die eigene Arbeit übernommen, besteht die \textbf{Pflicht}, diese zu kennzeichnen.
Wird dies unterlassen, liegt ein \textbf{Täuschungsversuch} (Plagiat) vor und die Arbeit kann als \textbf{nicht bestanden} bewertet werden.

Laut § 38 Abs. 2 der \emph{Allgemeinen Prüfungsbestimmungen} liegt \glqq\ [...] ein Täuschungsversuch [...] vor, wenn eine falsche Erklärung nach §§ 22 Abs. 7, 23 Abs. 7 abgegeben worden ist oder ein anderes Werk, eine Bearbeitung eines anderen Werkes, eine Umgestaltung eines anderen Werkes ganz oder teilweise in der Prüfungsarbeit wiedergeben werden, ohne dieses zu zitieren (Plagiat).\grqq~\cite{APB}


Die Angabe der Quellen in Ihrer Arbeit dient drei Zwecken, nämlich
\begin{itemize}
  \item der Kenntlichmachung nicht eigener Inhalte in Hinblick auf die Bewertung der Prüfungsleistung,\footnote{Wobei auch die Literaturarbeit an sich eine Leistung darstellt, die bewertet wird. Wenn Sie alles selber erarbeitet haben, dann ist dies in der Regel auch nicht die bestmögliche Leistung, da die Sichtung und das Auseinandersetzen mit den bestehenden Ansätzen, ein gegebenes Problem zu lösen, ebenfalls zu der erwarteten Leistung gehört.}
  \item der Nachvollziehbarkeit und Einordnung Ihrer Ergebnisse sowie
  \item der Anerkennung der/des ursprünglichen Urheberin/Urhebers des zitierten Gedankens oder Verfahrens.
\end{itemize}
Aus der Berücksichtigung dieser Ziele ergeben sich folgende Punkte, die zu beachten sind.

\paragraph{Zitierfähigkeit}
Idealerweise zitieren Sie veröffentlichte Werke aus allgemein zugänglichen Quellen (Bücher, Artikel, \etc).

Quellen, bei denen die Verfügbarkeit (der verwendeten Version) nicht garantiert werden kann (Internetquellen) oder der Urheber nicht klar nachvollziehbar ist (Wikipedia, o.\,Ä.), sollten möglichst vermieden werden.
Wird doch solch eine Quelle zitiert, ist die Version \bzw der Abrufzeitpunkt zu nennen und idealerweise die aktuelle Version zum Zeitpunkt des Zitates in digitaler Form der Arbeit beizulegen.

In Hinblick auf den erstgenannten Punkt ("`Kenntlichmachung nicht eigener Inhalte"') ist es auch möglich, "`persönliche Gespräche"' als Quelle anzugeben.
Dies sollte aber die Ausnahme sein, da hierbei der zweite Punkt ("`Nachvollziehbarkeit"') nicht gegeben ist.
In der Regel sollten solche Quellen nur bei der Problemanalyse bei praktischen, industrienahen Arbeiten vorkommen. (Interviews mit Anlagenfahrern, und ähnliches.)

Auch studentische Arbeiten sind zitierbar.
Wenn diese nicht in elektronischer Form vorliegt, ist es allerdings schwerer diese zu erhalten, was die Nachvollziehbarkeit erschwert.
Daher wäre es auch hier vorteilhaft, eine in diesem Sinne geeignetere Quelle zu finden.
Auch wenn theoretische Ergebnisse übernommen werden, deren Herleitung nicht in Ihrer Arbeit wiederholt wird, ist es günstiger eine Quelle zu nehmen, bei der von einer höheren Verbereitung (Lehrbuch) oder einem strengeren Review-Prozess (Veröffentlichungen in entsprechenden Zeitschriften) ausgegangen werden kann, da dies das "`Vertrauen"' in die Ergebnisse erhöht,\footnote{Was keinesfalls heißen soll, dass in Lehrbüchern oder Zeitschriftenveröffentlichungen keine Fehler sein können!} was auch der Nachvollziehbarkeit zuträglich ist.

Prinzipiell gilt bei der Übernahme nicht-eigener Gedanken und Ergebnisse, dass lieber eine im beschriebenen Sinne etwas ungeeignetere Quelle angegeben wird, als gar keine!
Letzteres ist ein Täuschungsversuch!


\paragraph{Wörtliche Zitate}
Wörtliche Zitate in ingenieurwissenschaftlichen Arbeiten sind unüblich.
Sollte doch ein wörtliches Zitat in die Arbeit übernommen werden, muss dieses buchstaben- und zeichengetreu, inklusive eventueller Rechtschreibfehler übernommen werden.
(Um deutlich zu machen, dass ein Fehler übernommen ist, können Sie diesen mit "`(sic!)"'\footnote{von lateinisch "`wirklich so"'} kennzeichnen.)
Das wörtliche Zitat wird in Anführungszeichen eingefasst.

Wenn Sie eine fremdsprachliches Zitat in einer übersetzten Form verwenden, dann wird dieses auch in Anführungszeichen mit einem folgenden Hinweis wie "`(Übersetzung aus dem \ldots vom Autor)"' geschrieben.


\paragraph{Sinngemäße Zitate}
Weit häufiger werden in wissenschaftlichen Arbeiten Ideen oder Meinungen anderer Autoren sinngemäß übernommen.
Diese müssen durch einen Verweis auf die Quelle gekennzeichnet werden.
Durch die Position des Verweises muss der Umfang der sinngemäßen Übernahme klar hervorgehen.

Steht die Quellenangabe innerhalb eines Satzes, so bezieht sich diese auf den Teil des Satzes vor dieser Quellenangabe.
Steht sie nach einem Satz, so wird diese in der Regel so verstanden, dass sie sich auf den aktuellen Absatz bis zu dieser Stelle bezieht.
Eine Quellenangabe sollte möglichst eng gefasst sein und sich nur in Ausnahmefällen auf größere Textabschnitte als einen Absatz beziehen.
Im letztgenannten Fall ist vor dem betreffenden Abschnitt klar zu machen, dass das Folgende auf der anzugebenden Quelle beruht.

Es ist dabei auch sehr wichtig, dass Sie den Inhalt der Quelle in \emph{eigenen Worten} wiedergeben.
Das gilt auch für Teilsätze!
Wenn Sie aus bestimmten Gründen eine zusammenhängende Folge von Worten übernehmen wollen, dann ist diese ein wörtliches Zitat und entsprechend zu kennzeichnen.

Eine Ausnahme hiervon stellen Gleichungen und mathematische Sätze dar, bei denen immer davon ausgegangen wird, dass diese wörtlich (abgesehen von Anpassungen der Notation wie Variablennamen) übernommen sind, wenn eine Quelle angegeben ist.
Sollten Sie dabei Umformulierungen vornehmen, um diese beispielsweise zu vereinfachen, sollte klar angegeben werden, welche dies sind.
(\ZB "`nach [..., Satz 3.2], dort allgemein für komplexe Zahlen formuliert, hier speziell auf reele Zahlen angewandt"'.)

Wenn es Sie beim Schreiben "`nervt"', dass Sie absatzlang damit zu kämpfen haben, eigene Formulierungen für den Inhalt einer Quelle zu finden, dann liegt es meist daran, dass Sie wahrscheinlich etwas in einer Ausführlichkeit zitieren, die weder gefordert noch gewollt ist.
Sie müssen prinzipiell keine Herleitungen aus anderen Quellen in Ihrer Arbeit wiederholen.
Hier reicht es in der Regel, bei der Verwendung der Ergebnisse auf die entsprechenden Quellen zu verweisen.
Eine Ausnahme läge beispielsweise vor, wenn Sie später die Eigenschaften eines Verfahrens disktuieren wollen, und hierzu vorteilhafterweise auf die Herleitung bezug nehmen.
Oder wenn Sie eine Herleitung verwenden, die nur sehr schwer zu finden ist, für das Verständnis Ihrer Arbeit aber eine wesentliche Rolle spielt.
(Hierfür würde sich dann \zB auch der Anhang anbieten. Auf jeden Fall ist dies dann klar zu kennzeichnen und darauf zu achten, keinesfalls wörtlich zu übernehmen.)


\paragraph{Grundwissen}

In jedem Zweig der Wissenschaft gibt es ein bestimmtes Grundwissen, das nicht mit Quellen hinterlegt werden muss, da allen, die in dem Bereich tätig sind, klar ist, dass diese Ergebnisse erstens nicht von Ihnen stammen und zweitens auch stimmen.
Auch sorgt dies für einen besser lesbaren Text, als wenn man bei strengster Auslegung auch eine Quelle für die Anwendung der Grundrechenarten angeben würde.

Es existiert keine Checkliste von Begriffen oder Gleichungen, die ohne Quellenangabe verwendet werden können, zumal dies auch von der entsprechenden Fachrichtung abhängt.

Zur Hilfestellung bei der Abschätzung, ob eine bestimmte Aussage diesem Grundwissen zuzurechnen ist, sei hier auf eine entsprechende Internetseite%
\footnote{Academic Integrity at MIT. A Handbook for Students. What is Common Knowledge?\\
URL: https://integrity.mit.edu/handbook/citing-your-sources/what-common-knowledge (abgerufen am 14.01.2021)}
(in englischer Sprache) des MIT verwiesen, die zur Beantwortung dieser Fragestellung drei Testfragen angibt und ein paar Beispiele dazu liefert. Die erste dort aufgeführte Testfrage zielt darauf ab, sich klarzumachen, wer der Adressat Ihrer Arbeit ist. Basierend auf dieser Zielgruppe werden auf der zitierten Seite dann die beiden folgenden Fragen gestellt.
\begin{itemize}
  \item Von welchem Wissensstand kann ich bei dieser Gruppe ausgehen?
  \item Würde ich aus dieser Gruppe gefragt werden, woher die betrachtete Aussage stammt?
\end{itemize}
Für die Anwendung dieser Methode auf Ihre Arbeit sollen Sie bei der Zielgruppe hier nicht an den Betreuer denken, sondern an Studierende der gleichen Fachrichtung, die sich im Studium ungefähr an der gleichen Stelle befinden, und nicht unbedingt die gleichen Vertiefungsfächer gehört haben. Auch ist zu bedenken, dass auch in den Grundlagenfächern durchaus je nach Professor der eine oder andere weitergehende Aspekt behandelt wird, der nicht jedem Bachelor- oder Masterstudent bekannt sein wird.

Prinzipiell zeigen diese Überlegungen, dass die Frage nach dem Grundwissen nicht nur die Fachrichtung, sondern auch die Art der Arbeit mit einschließt. Es lässt sich also keine einfache Liste angeben. Daher sind im Folgenden lediglich ein paar Beispiele genannt, bei denen wir bei Bachelor- und Masterarbeiten in unserem Bereich (noch) keine Quellenangabe erwarten, und ein paar Beispiele für den Fall, bei denen dies (eher) der Fall ist.
Damit soll diese Grenze angedeutet werden.
Bei Unsicherheit können Sie mit Ihrem Betreuer Rücksprache halten, und im Zweifel lieber eine Quellenangabe mehr als zu wenig verwenden.

Für die elementaren mathematischen Begriffe müssen in ingenieurswissenschaftlichen Arbeiten keine Quellen genannt werden.
So können Sie von natürlichen, rationalen, reellen und auch komplexen Zahlen reden, ohne hierfür eine Quelle anzugeben. Auch die "`elementaren"' Techniken der Integration und Differentiation bedürfen keines Verweises.
Der Satz des Pythagoras sowie die Definitionen von Sinus- und Kosinus können ebenfalls ohne Quelle verwendet werden.
Wenn Sie bestimmte Integrale in einem Tabellenwerk nachschlagen, ist dies natürlich anzugeben. Ebenso sollte bei der Verwendung der Additionstheoreme auf entsprechende Stellen verwiesen werden, an denen diese nachgeschlagen werden können.

Die Grundgleichungen der elementaren, linearen elektrotechnischen Elemente (Widerstand, Spule, Kondensator) können sie ebenso verwenden, wie die der entsprechenden mechanischen Elemente.
Bei den Newton'schen Axiomen, die Sie möglicherweise zum Aufstellen der Systemgleichungen verwenden wollen, reicht der Hinweis auf den Namen.
Wenn Sie aber \zB den Lagrang'schen Formalismus verwenden, sollte ein Hinweis auf eine Quelle angegeben werden, in der dieser erläutert wird.

Bei einer regelungstechnischen Arbeit wird nicht erwartet, dass Sie den Aufbau und die Funktion eines Standardregelkreises mit einer Quelle hinterlegen.
Jedoch sollten Sie schon möglicherweise als "`gering"' empfundene Variationen, wie eine Zwei-Freiheitsgrade-Struktur, belegen.
Sie müssen den Begriff einer Übertragungsfunktion oder einer Zustandsraumdarstellung nicht mit einer Quelle hinterlegen.
Sämtliche Reglerauslegungsverfahren sollten bei Nennung jedoch belegt werden.


\paragraph{Quellenangaben im Literaturverzeichnis}
Für die Darstellung der Verweise, als auch für die Darstellung der Quellen im Literaturverzeichnis gibt es verschiedene Zitierweisen.
Üblich sind die Harvard-Variante mit Autor und Veröffentlichungsjahr in runden Klammern, wie zum Beispiel (Isermann, 2001) oder eine fortlaufende Nummerierung in eckigen Klammern, wie in dieser Vorlage.

Die Quellenangabe im Literaturverzeichnis muss alle notwendigen Informationen enthalten, um die genannte Quelle auf"|finden zu können.
Die genauen Angaben hängen von der Art der Quelle ab.
Zur Orientierung sind hier die Angaben für die wichtigsten Arten genannt:
\begin{itemize}
  \item Artikel:\\
    Autor, Titel, Zeitschrift, Jahr, Ausgabe, Nummer, Seite, \ggf DOI
  \item Buch:\\
    Autor oder Herausgeber, Titel, Verlag, ISBN, Jahr\\
    (Der Verlagsort, der früher immer angegeben wurde, spielt durch die ISBN eigentlich keine große Rolle mehr, die Quelle zu finden.
    Wenn Sie diesen angeben, reicht aber auch bei größeren Verlagen die Angabe eines Ortes aus.)
  \item Dissertation:\\
    Wenn als Buch herausgegeben, dann als Buch zitieren, \ggf mit dem Hinweis "`Dissertation"'.\\
    Sonst: Autor, Titel, \emph{Dissertation}, Hochschule, Jahr
  \item Skript (wenn nicht als Buch mit ISBN herausgegeben), studentische Arbeit:\\
    Autor, Titel, Art (Vorlesungsskript, Masterarbeit, \ldots), Hochschule, Jahr
  \item Internetquelle:\\
    Autor (wenn erkennbar), Seite mit Link (idealerweise auf feste Version), Datum und Uhrzeit des Abrufs
  \item Sonstiges (Handbücher, technische Berichte, Normen):\\
    Zu den Angaben sich sinngemäß an den obigen Beispielen orientieren.
    Technische Unterlagen haben meist eine Revisionsnummer, die angegeben werden kann.
\end{itemize}



\paragraph{Quellenverweise im Text}

Bei den Quellenverweisen im Text sollte -- insbesondere bei umfangreicheren Werken -- auch eine Seitenangabe angegeben sein.
Bei mathematischen Sätzen oder Gleichungen kann auch auf die entsprechende Nummerierung der Quelle zurückgegriffen werden.
Dies dient der Nachvollziehbarkeit Ihrer Arbeit.
(Und hilft auch Ihnen, zu einem späteren Zeitpunkt die zitierte Stelle wieder zu finden.)




\subsubsection*{\LaTeX}

Das Paket \texttt{biblatex} ermöglicht die automatisierte Erzeugung eines Literaturverzeichnisses, wobei zwischen verschiedenen Zitierweisen gewählt werden kann.

Zunächst benötigt man eine Literaturliste, die alle Quellen, die zitiert werden sollen, enthält.
Dies ist eine Datei mit der Endung \verb|bib|, die die verschiedenen Quellen in einem bestimmten Format enthält.
Diese kann man \zB mit dem freien Programm \textsc{JabRef} erzeugen.
Es gibt aber auch viele weitere Programme, die einen größeren Funktionsumfang zur Literaturverwaltung bieten (\textsc{JabRef} ist im Grunde "`lediglich"' ein Viewer für \verb|bib|-Dateien) und die ihre Datenbanken auch in das bibtex- \bzw biblatex-Format exportieren können.

Für die Erstellung des Literaturverzeichnisses aus der Literaturliste bietet sich das Paket \verb|biblatex| zusammen mit dem zusätzlichen Programm \textsc{Biber} an.\footnote{Dies ist der im Minimalbeispiel voreingestellte Weg.}
Die Sortierung des Literaturverzeichnisses kann alphabetisch (voreingestellt, oder \bspw durch die Paketoption \texttt{sorting=nyt}) oder nach dem Erscheinen der Verweise erfolgen (durch die Paketoption \texttt{sorting=none}).

Ein Verweis wird im Text an der entsprechenden Stelle mit \verb|\cite{<label>}| eingefügt, wobei \verb|<label>| der entsprechende Bezeichner des Eintrages der Literaturliste ist.
Schließt der Verweis einen Satz ab, folgt der Punkt \emph{hinter} dem Verweis.
Zusätzliche Angaben wie Seiten(bereiche) oder eine Kapitelnummer können in eckigen Klammern als optionaler Parameter des Befehles \verb|cite| angegeben werden, \zB \verb|\cite[123f]{Isermann2001}|.



\section{Rechtschreibung}

Studentische Abschlussarbeiten am IAT sind nach den aktuell geltenden Regeln der deutschen Rechtschreibung zu verfassen, \cite{Duden}.
Im Internet findet man unter \url{http://www.duden.de/} einen Crashkurs zur deutschen Rechtschreibung.

Wird die Arbeit in Englisch abgefasst, ist es freigestellt, ob die amerikanischen oder britischen Schreibweisen verwendet werden.
Es ist jedoch darauf zu achten, über den gesamten Text einheitlich eine der beiden Varianten zu verwenden.


\subsubsection*{\LaTeX}

Das Paket \verb|babel| kann über die Option \verb|main=ngerman| für die "`neue"' deutsche Rechtschreibung geladen werden.
Es ist aber zu beachten, dass \verb|babel| kein Paket zur Überprüfung der Rechtschreibung ist, sondern dazu dient, sprachabhängige Bezeichnungen (\zB "`Inhaltsverzeichnis"' anstelle von "`Table of Contents"') zu laden sowie die sprachspezifischen Trennregeln anzuwenden.
Letzteres ist auch der Grund, dass man zusätzlich zur Hauptsprache \verb|ngerman| auch \verb|english| laden sollte, wenn man eine Kurzzusammenfassung auch in Englisch schreiben möchte.
Es ist dann für den Text dieser Zusammenfassung die Sprache nach \verb|english| zu wechseln.

Bei einer englischen Abfassung der Arbeit ist entsprechend die Hauptsprache \verb|english| (amerikanische Regeln) oder \verb|british| zu laden.

Insbesondere wenn man nicht besonders sicher in der Sprache, in der die Arbeit abgefasst wird, ist, kann es sich auch lohnen, den Inhalt der Arbeit am Ende einmal in Microsoft Word zu kopieren.
Es sind dann natürlich noch entsprechende Anpassungen an dem Text vorzunehmen, da Word möglicherweise verwendete \LaTeX-Befehle nicht kennt, dafür verfügt es über eine natürlich nicht fehlerfreie, aber erfahrungsgemäß sehr leistungsfähige Überprüfung von Rechtschreibung und Grammatik.




\section{Gliederung des Dokuments}
\label{sec:Latex-Gliederung}
Im Inhaltsverzeichnis wird die Gliederung der Arbeit dargestellt.
Die Überschriften und Seitenangaben der Kapitel, Unterkapitel und weiteren Abschnitte müssen mit den Elementen im Inhaltsverzeichnis übereinstimmen.
Überschriften sind kurz und prägnant zu formulieren und dürfen keine vollständigen Sätze sein.
Gibt es Unterpunkte in der Gliederung, so müssen immer mindestens zwei davon existieren und inhaltlich auf der gleichen Ebene sein.
Die einzelnen Punkte des Inhaltsverzeichnis müssen nummeriert werden.
Die Übersichtlichkeit des Inhaltsverzeichnisses kann durch Einrücken der Unterpunkte erhöht werden.

Absätze gliedern ebenfalls den Text und sollten entsprechend verwendet werden.


\subsubsection*{\LaTeX}

Das Inhaltsverzeichnis wird in \LaTeX\ automatisch erstellt.
Die Strukturierung des Dokumentes wird dazu mit entsprechenden Befehlen vorgenommen.
Innerhalb der einzelnen Kapitel \verb|\chapter{...}| werden weitere Unterteilungen mit den Befehlen \verb|\section{...}|, \verb|\subsection{...}| \usw vorgenommen.
Werden diese mit einem \verb|*| versehen, dann erhält der jeweilige Abschnitt keine Nummer und erscheint nicht im Inhaltsverzeichnis.
Dies kann in manchen Fällen nützlich sein.

Um eine korrekte Darstellung des Inhaltsverzeichnisses zu erhalten, muss \ggf mehrmals hintereinander kompiliert werden, da sich aufgrund von Gleitobjekten Seitenzahlen ändern können.
Dreimaliges Kompilieren reicht in der Regel.


Werden Gleichungen oder Listen in einer Umgebung in den laufenden Text eingefügt, darf dazwischen kein Absatz (\dah eine Leerzeile im Quelltext) stehen.%
\footnote{Dies führt dazu, dass sich   damit unnatürlich große Zwischenräume ergeben.
  Wenn nach der Gleichung oder der Liste ein neuer Sinnzusammenhang beginnt, wäre dies gewollt.
  Wenn der eigentliche Absatz aber durch die Gleichung oder Liste hindurchläuft, ist dies unerwünscht.}
Um den Quelltext besser zu gliedern, kann an dieser Stelle eine Zeile mit dem Kommentarzeichen \verb|%| eingefügt werden.
Eine Leerzeile darf nur dann im Quelltext stehen, wenn auch wirklich ein Absatz erwünscht ist.
Der Zeilentrenner \verb|\\| erzeugt übrigens keinen Absatz und darf im laufenden Text überhaupt nicht vorkommen.




\section{Gleitobjekte und Referenzierungen}

Aus Gründen eines gefälligen Textsatzes werden Bilder und Tabellen (sowie \zT auch Codeausschnitte und ähnliches) nicht zwangsweise an die Stelle gesetzt, an der diese im Text zum ersten Mal erwähnt werden, sondern meist an den Anfang oder das Ende einer Seite.
Daher erhalten Bilder und Tabellen entsprechende Unter- \bzw Überschriften mit einer Nummerierung und einer kurzen Beschreibung (\zB "`Abbildung 12: Standardregelkreis"' oder "`Bild 12: Standardregelkreis"').

Im Text werden dann die verschiedenen Objekte über die entsprechende Bezeichnung und Nummer referenziert, also \zB "`\ldots wie in Abbildung~12 erkannt werden kann \ldots"'.

Genauso kann auf Kapitel, Abschnitte, (mathematischen) Sätze, Gleichungen und andere "`Objekte"' referenziert werden.
Gleichungen nehmen hierbei eine Sonderrolle ein, da diese über eine Zahl in runden Klammern referenziert werden ("`Wie man aus Gl. (2.2) erkennen kann, \ldots"').

Es ist darauf zu achten, ob es sich um \emph{Kapitel} oder \emph{Abschnitte} handelt.
(Also, "`\ldots wie schon in Kapitel~3 diskutiert wurde \ldots"', aber "`\ldots wie schon in Abschnitt~3.2 diskutiert wurde \ldots"'.)

Die Verweise sind vollständig, inklusive der Bezeichung der Art des "`Objektes"', zu bezeichnen.
\Dah es kann nicht "`\ldots, siehe 4.2."' geschrieben werden, sondern es muss \zB "`\ldots, siehe Abbildung~4.2."' oder "`\ldots, siehe Abschnitt~4.2."' geschrieben werden, um den Verweis eindeutig zu machen.
Eine Ausnahme stellen hierbei Gleichungen dar.
Aufgrund der anzuwendenen Konvention, die Gleichungsnummern in Klammern zu schreiben, ist die Referenzierung über diese Nummer alleine eindeutig, und es kann (muss aber nicht) "`\ldots, siehe (4.2)."' geschrieben werden.


\subsubsection*{\LaTeX}

Bilder und Tabellen werden in der Regel in eine \verb|figure|- \bzw \verb|table|-Umgebung eingeschlossen.
Damit werden diese zu sogenannten \emph{Gleitobjekten}, \dah sie erscheinen nicht an der Stelle, an der sie eingebunden werden, sondern oben oder unten auf einer Seite.
Optional können in \verb|[]| noch Positionierungswünsche angegeben werden.
Diese wären \verb|h| (here), \verb|t| (top), \verb|b| (bottom) und \verb|p| (page).
Letzteres bedeutet, dass man \LaTeX{} erlaubt, eine Seite zu erzeugen, die nur Gleitobjekte enthält.
Dies ist zu verwenden, wenn entweder sehr große Abbildungen oder Tabellen, die eine ganze Seite einnehmen, oder sehr viele Abbildungen hintereinander gezeigt werden sollen.
Erlaubt man nicht explizit die Plazierung auf einer eigenen Seite, werden diese dann bis an das Ende des Dokumentes verschoben.
In dem Kontext kann auch der Befehl \verb|\FloatBarrier| aus dem Paket \verb|placeins| nützlich sein.
Dieser zwingt \LaTeX{} alle noch nicht plazierten Gleitobjekte an der genannten Stelle abzulegen, bevor der weitere Text gesetzt wird.
(Man sollte dies auch nicht übermäßig verwenden. Bei Arbeiten in unserem Gebiet kann es aber schon einmal vorkommen, dass ein paar zusammengehörende Ergebnisplots gezeigt werden sollen.)

Mit dem Befehl \verb|\caption{...}| erhalten Bilder eine \emph{Unterschrift} und Tabellen eine \emph{Überschrift}.
Dazu ist bei Bildern die \verb|caption| \emph{nach} dem Einfügen des Bildes innerhalb der \texttt{figure}-Umgebung (siehe unten) und bei Tabellen die \verb|caption| \emph{vor} der inneren \verb|tabular|-Umgebung zu setzen.

Alles, was sinnvoll referenziert werden kann (Kapitel, Abschnitte, Bilder, Tabellen, Gleichungen, mathematische Sätze, \ldots), kann in \LaTeX{} mittels \verb|\label{...}| benannt werden.
Dadurch ist es möglich, sie später mit \verb|\ref{...}| (\bzw \verb|eqref{...}| für Gleichungen) oder \verb|\pageref{...}| zu referenzieren.

Es empfiehlt sich, den Namen (Labels) eine Markierung voranzustellen, aus der hervorgeht, um welches Objekt es sich handelt.
Üblich sind \verb|cha:| für \glqq Chapter\grqq, \verb|sec:| für \glqq Section\grqq, \verb|fig:| für \glqq Figure\grqq, \verb|tab:| für \glqq Table\grqq\ und \verb|eq:| für \glqq Equation\grqq.
Ein Bild benennt man also \zB mit \verb|\label{fig:Ausgangssignal}|.

Vor \verb|\ref{...}| sollte ein \emph{festes} Leerzeichen ("`\verb|~|"') stehen, damit dort kein Umbruch erfolgen kann.




\section{Bilder}
\label{sec:Latex-Bilder}

Wenn möglich, sollten Bilder als Vektorgrafik eingebunden werden, damit sichergestellt wird, dass alle Details beim Ausdrucken erhalten bleiben.
Es ist dabei auf eine ausreichende Strichstärke zu achten.
Ebenfalls sollte die im Bild verwendete Schrift die gleiche sein, wie im übrigen Dokument.

Werden doch Pixelgrafiken verwendet, so ist die richtige Wahl der Auf"|lösung von besonderer Bedeutung.
Einerseits sollten die Bilder auf dem ausgedruckten Dokument gut aussehen, andererseits aber auch eine zügige Bildschirmdarstellung und kleine Dateigröße ermöglichen.

In \appref{cha:Anhang-Grafiken} sind geeignete Programme zur Erstellung von Bildern und Plots mit ihren Eigenschaften aufgelistet.

Generell ist die Arbeit (und insbesondere die Grafiken) so zu gestalten, dass sie auch schwarzweiß gedruckt werden kann.
Farbige Fotos und Screenshots verursachen dabei \iA keine zusätzlichen Probleme.
Werden jedoch \zB farbige Kurven in einem Diagramm dargestellt, hat dies folgende Konsequenzen:
\begin{itemize}
  \item Keine zu hellen Farben für die Linien verwenden, da diese sonst beim Drucken nicht zu erkennen sind.
    RGB-Grün $(0,1,0)$ ist tabu!
  \item Bei mehreren Kurvenverläufen darf deren Farbe nicht das einzige Unterscheidungsmerkmal sein: Entweder sind unterschiedliche Linienformen zu verwenden oder die Kurven im Diagramm beschriften.
\end{itemize}

Bei der Erstellung von Plots kann sich an der \href{http://de.wikipedia.org/wiki/DIN_461}{\texttt{DIN 461}} (Grafische Darstellung in Koordinatensystemen) orientiert werden.
Hilfreich in diesem Zusammenhang und überhaupt für die korrekte Schreibweise von Zahlen und Einheiten im Fließtext und in Formeln sind die Hinweise der TU-Chemnitz: \url{www.tu-chemnitz.de/physik/FPRAK/Grundsatz/Literatur/si_v1.pdf}.



\subsubsection*{\LaTeX}

Eine Grafik lässt sich am einfachsten mit dem Befehl \verb|\includegraphics{}| aus dem \verb|graphicx|-Paket einbinden.
In welchem Format \verb|graphicx| die Grafiken benötigt, hängt davon ab, welches \LaTeX-Programm man verwendet.
Bei der Verwendung von \verb|pdflatex| und neueren Programmen sollten Grafiken im \emph{Portable-Document-Format} (\verb|*.pdf|) sowie die Pixelformate \texttt{jpeg} und \texttt{png} direkt eingebunden werden können.

Bietet das verwendete Programm zur Erzeugung der Grafik keinen direkten Export von PDF-Grafiken an, so lässt sich möglicherweise aus einer im EPS-Format exportierten Grafik mit Hilfe des Acrobat Distiller oder Ghostscript ein PDF erzeugen.

Bilder müssen zentriert sein (\verb|\centering|) und eine Bild\emph{unterschrift} (\verb|\caption{...}|) besitzen.
Um auf eine Abbildung zu referenzieren, muss sie mit einem \verb|\label{fig:...}| versehen werden.
Nur in besonderen Ausnahmefällen sollten Bilder mit Text umflossen werden.

Wurden Abbildungen einer Quelle entnommen, muss dies entsprechend mit einem Verweis auf die Quelle im Literaturverzeichnis gekennzeichnet werden.
%Wenn dazu das Makro \verb|\cite| verwendet wird, so ist zusätzlich das optionale Argument von \verb|\caption| ohne den Literaturverweis zu verwenden, damit die Verlinkungen im Abbildungsverzeichnis nicht zerstört werden.
Werden Abbildungen einer Quelle nachempfunden, angepasst oder abgeändert, ist dies mit dem Zusatz \glqq in Anlehnung an [...]\grqq\ oder ähnlich anzugeben.

\begin{figure}[htp]
  \centering
  \input{./tikz/BSB_Beispiel.tikz}
  \caption{Standard-Regelkreis; Bild erstellt mit \texttt{TikZ}}
  \label{fig:Standardregelkreis}
\end{figure}

Zum Erstellen von Bildern bietet es sich an, die \LaTeX-Pakete \texttt{tikz} (für Blockschaltbilder und andere Zeichnungen) und \texttt{pgfplots} (für Plots) zu verwenden.
Der große Vorteil davon liegt darin, dass die Grafiken direkt in \LaTeX{} kompiliert werden und somit die gleiche Schriftart- und -größe besitzen und die Liniendicke definiert ist.
Dadurch ergibt sich ein schönes Gesamtbild des Dokuments, das wirkt, als wäre es aus einem Guss, siehe \zB \figref{fig:Standardregelkreis}.

Für Matlab existiert die Funktion \verb|matlab2tikz| (über den Mathworks File Exchange herunterladbar), welche Matlab-Plots in das \texttt{pgfplots}-Format konvertiert.
Hierzu ist auch die dazugehörende Funktion \verb|cleanfigure| zu erwähnen, welches die Plotpunkte reduziert, was zu einer Reduzierung der "`Kompilierungszeit"' und der Dateigröße führt.

\paragraph*{\textsc{tikz}-Externalisierung}
Durch das Kompilieren des Bildes direkt in \LaTeX{} kann sich jedoch auch ein Problem ergeben, nämlich dann, wenn die Datenmenge zu groß ist.
Der Kompiler beschwert sich dann mit einem\\
\texttt{[...] ! TeX capacity exceeded, sorry [main memory size=3000000]}\\
oder ähnlichem.
Dies wird sehr schnell erreicht, wenn zum Beispiel viele Plots in einem Dokument vorhanden sind.

Abhilfe dagegen schafft das separate Kompilieren der Bilder und anschließende Einbinden derselben als pdf-Dateien.
Dazu ist \zB möglich, den in \LaTeX{} eingebauten Mechanismus des Paketes \texttt{tikz} zu verwenden, indem die \texttt{external} Bibliothek durch \verb|\usetikzlibrary{external}| eingebunden wird, die sich um alles Weitere kümmert.

Damit die in \texttt{tikz} eingebaute Externalisierung funktioniert, muss das Aufrufen von Systembefehlen durch \LaTeX{} erlaubt sein, was sich je nach verwendeter Distribution durch das Hinzufügen der Kommandozeilenparameter \texttt{-shell-escape} oder \texttt{-enable-write18} zum Kompilierbefehl erreichen lässt.
Unter \Texstudio{} kann dazu unter \zitat{Befehle} bei pdf\LaTeX{} und \LaTeX{} die entsprechende Zeile geändert werden.



\section{Tabellen}
Tabellen müssen ebenfalls zentriert sein und besitzen eine zentrierte Tabellen\emph{überschrift}. Hier gelten die gleichen Regeln zur Quellenangabe wie bei den Bildern.

Damit Tabellen \glqq schön\grqq\ aussehen, empfiehlt es sich, einige Grundregeln zu beachten.
Es gilt das Prinzip: \emph{weniger ist mehr}.
So sollte auf die Verwendung von senkrechten Linien verzichtet werden und nur wichtige Zeilen, wie zum Beispiel Überschriften, Sinnabschnitte, Unterpunkte, \etc mit horizontalen Linien getrennt werden.

\begin{table}[htb]
  \begin{center}
    \caption{Parameter}
    \label{tab:Systemparameter}
    \begin{tabular}{ccclp{3mm}cccl}
      \toprule
          &  \multicolumn{1}{c}{Ref.}  &  \multicolumn{1}{c}{Mod.}  &  Einheit          &  &          &  \multicolumn{1}{c}{Ref.}  &  \multicolumn{1}{c}{Mod.}  &  Einheit\\\cmidrule(r){1-4}\cmidrule(l){6-9}
      $m_1$  &  4,0    &  4,63  &  \unit{kg}        &  &  $d_1$      &  0    &  0    &  \unit{\frac{N\,m\,s}{rad}}\\[1mm]
      $m_2$  &  10,1  &  11,15  &  \unit{kg}        &  &  $d_2$      &  0    &  0    &  \unit{\frac{N\,m\,s}{rad}}\\[1mm]
      $m_3$  &  45,7  &  42,5  &  \unit{kg}        &  &  $l_1$      &  0,5    &  0,45  &  \unit{m}\\[1mm]
      $J_1$  &  0,967  &  0,993  &  \unit{kg\,m^2}      &  &  $l_2$      &  1,5    &  1,59  &  \unit{m}\\[1mm]
      $J_2$  &  0,571  &  0,599  &  \unit{kg\,m^2}      &  &  $\varphi_{1,0}$  &  100    &  98,5  &  \unit{\degree}\\[1mm]
      $g$    &  9,81  &  9,81  &  \unit{\frac{m}{s^2}}  &  &  $\varphi_{2,0}$  &  5    &  4,66  &  \unit{\degree}\\
      \bottomrule
    \end{tabular}
  \end{center}
\end{table}


\subsubsection*{\LaTeX}

(Kleinere) Tabellen, wie \zB \tabref{tab:Systemparameter}, werden in einer \verb|tabular|-Umgebung gesetzt, die in der Regel in einer \verb|table|-Umgebung eingebettet ist.
Für größere Tabellen, die über eine Seite hinausgehen, gibt es besondere Umgebungen.
Es existieren viele vordefinierte Spaltentypen und es können auch eigene Typen definiert werden, um Zellen mit bestimmten Eigenschaften (horizontale und vertikale Ausrichtung, feste oder variable Breite) zu erhalten.
Auch können einzelne Gruppen von Zellen horizontal oder vertikal verbunden werden.
Das Paket \verb|booktabs| stellt gesonderte Linientypen für den Kopf und Fuß einer Tabelle zur Verfügung.

Ein Referenzieren wird auch hier mit \verb|\label{tab:...}| ermöglicht.




\section{Mathematische Formeln}

Mathematische Formeln, wie
\begin{equation}
  \label{eq:Allgemein-Approx}
  \int\limits_0^\infty g(x)\ud{}x \approx \sum_{i=1}^n w_i\eexp{x_i} g(x_i)\;,
\end{equation}
werden eingerückt linksbündig dargestellt und nur dann nummeriert, wenn auf sie im Text verwiesen wird.
Für eine bessere Lesbarkeit ist es sinnvoll, Gleichungen als Bestandteil des Satzes aufzufassen, der sich auf natürliche Weise "`mitlesen"' lassen sollte, wie \zB bei \equref{eq:Allgemein-Approx} demonstriert.
Die dabei möglicherweise nötigen Satzzeichen sind mit etwas Abstand zu setzen.
Die Referenzierung einer erst später aufgeführten Gleichung sollte vermieden werden.

Gelingt die "`schöne"' Einbettung der Gleichung in den Fließtext nicht, so sollte diese zumindest nach einem Doppelpunkt stehen.
Keinesfalls sollte eine Gleichung "`motivationslos"' zwischen zwei Absätzen stehen.

In der Regel wird für mathematische Ausdrücke eine andere Schriftart als die des normalen Textes verwendet, wobei diese aufeinander abgestimmt sein sollten.
Es ist darauf zu achten, dass \emph{alle} mathematischen Ausdrücke (auch wenn es nur einzelne Zeichen sind) in der enstprechenden gleichen Schrift für mathematische Ausdrücke geschrieben werden.

Es gilt die Regel, dass gewöhnliche mathematische Größen \emph{kursiv} geschrieben werden, Ausdrücke mit konventioneller (feststehender) Bedeutung dagegen in normaler (aufrechter) Schrift, siehe~\cite{DIN1338}.
Es sind insbesondere Einheiten, Standardfunktionen und -operatoren sowie mathematische Konstanten aufrecht zu schreiben,
\begin{equation*}
  \eexp{ax}\;,\qquad a+\iu{}b\;,\qquad
  \int\limits_{t=0}^{t=1\,\unit{s}} f(t)\cdot\sin(\omega t)\ud{}t\;.
\end{equation*}
Konsequenter Weise muss dieses Prinzip auch auf Indizes angewendet werden, also \zB
\begin{equation*}
  \sum_{i,\,j}a_{ij}\,\sin(ijx)\;,\qquad\text{aber:}\qquad
  K_\mathrm{Regler} = 0,5\;.
\end{equation*}
Als Ausnahme von der oben genannten Regel werden große griechische Buchstaben meist \emph{nicht} kursiv geschrieben.

Matrizen und Vektoren können in \textbf{fetten} Buchstaben gesetzt werden.
Damit sie sich besser von den übrigen Symbolen abheben, werden dann aber nicht \textbf{\emph{fett-kursiv}} sondern \textbf{fett-aufrecht} geschrieben, also \zB
\begin{equation*}
  \mat{A} = \begin{bmatrix}
    a_{11} & \ldots & a_{1n}\\
    \vdots & \ddots & \vdots\\
    a_{n1} & \ldots & a_{nn}
  \end{bmatrix}\;,\qquad
  \ve{x} = \begin{bmatrix}
    x_1\\
    x_2
  \end{bmatrix}\;,\qquad
  \ve{\upbeta}^\transp = \begin{bmatrix}\beta_1  &  \ldots  &  \beta_m\end{bmatrix}
\end{equation*}

Chemische Formelzeichen schreibt man grundsätzlich in aufrechter Schrift.
Variable Größen sind aber auch hier kursiv:
\begin{equation*}
  \mathrm{H_2 O}\;,\qquad
  \mathrm{NO}_x\;,\qquad
  \mathrm{Fe_2^{2+}Cr_2^{\vphantom{2+}}O_4^{\vphantom{2+}}}
\end{equation*}


Beim Referenzieren von Gleichungen muss diese nummeriert werden. Dabei steht die Nummer der Gleichung immer in runden Klammern, wie bei Gl.~\eqref{eq:Allgemein-Approx}.

Alle Variablen und besondere Formelzeichen, die in einer Gleichung vorkommen, müssen entweder vor dieser Gleichung schon eingeführt sein oder spätestens unmittelbar nach dieser Gleichung erklärt werden.


\subsubsection*{\LaTeX}

Abgesetzte Gleichungen werden in einer entsprechenden mathematischen Umgebung erstellt.
Ob eine Gleichung nummeriert wird oder nicht, wird durch die Wahl der Umgebung, in der diese gesetzt wird, festgelegt.
Es ist empfehlenswert, die Gleichungsumgebungen des \verb|amsmath|-Paketes zu verwenden, da diese über eine Vielzahl an Möglichkeiten der Ausrichtung mehrerer oder längerer Gleichungen verfügen und dabei insgesamt zu einem einheitlichen Bild gelangen.
(Verwendet man die Gleichungsumgebungen, die in \LaTeX{} fest dabei sind, so passt der Abstand zum Gleichheitszeichen bei den untereinander ausgerichteten Gleichungen nicht zu den normalen Gleichungen.)
Mit der Umgebung
\begin{verbatim}
\begin{align}
    ...
\end{align}
\end{verbatim}
erzeugt man eine oder mehrere nummerierte Gleichungen (die zur Referenzierung dann noch innerhalb der \verb|align|-Umgebung mit einem \verb|label| zu versehen sind), während mit
\begin{verbatim}
\begin{align*}
    ...
\end{align*}
\end{verbatim}
nicht-nummerierte Gleichungen erstellt werden.
Innerhalb dieser Umgebungen dürfen keine Leerzeilen im Quelltext vorkommen. Um den Quelltext dennoch nach den eigenen Vorstellungen zu strukturieren, können "`beinahe-Leerzeilen"' verwendet werden, die lediglich das Kommentarzeichen \verb|%| enthalten.


Gleichungen oder mathematische Ausdrücke innerhalb des normalen Textes, für die kein Umbruch erzeugt werden soll, werden im mathematischen Modus \verb|$...$| geschrieben.
Letzteres ist auch dann konsequent zu machen, wenn es sich bei dem mathematischen Ausdruck um ein einzelnes Zeichen handelt, damit dieses in der richtigen Schrift erscheint.

Da in mathematischen Umgebungen alle Zeichen (mit Ausnahme der großen griechischen Buchstben) zunächst \emph{kursiv} gesetzt werden, müssen Ausdrücke, die in aufrechter Schrift erscheinen sollen, mit dem Befehl \verb|\mathrm{...}| gekennzeichnet werden.
Für aufrecht zu schreibende Standardfunktionen werden von \LaTeX\ die entsprechenden Befehle \verb|\sin|, \verb|\log|, \verb|\max| \usw zur Verfügung gestellt.
Soll innerhalb einer Formel "`richtiger"' Text erscheinen, ist \verb|\text{...}| anstelle von \verb|\mathrm{...}| zu verwenden.

Wenn nach abgesetzten Gleichungen ein Satzzeichen, also ein Komma oder ein Punkt, folgt, dann kann dieses nach einem (kleinen) Leerzeichen \verb|\,| gesetzt werden, um es optisch besser von der Gleichung zu trennen.
Den Abstand können Sie nach Geschmack festlegen, es muss aber einheitlich sein.



\section{Auszeichnungen und Hervorhebungen}
Wichtige Begriffe werden durch eine andere Schrift hervorgehoben (ausgezeichnet).
Man unterscheidet dabei integrierte und aktive Auszeichnungen.
Integrierte Auszeichnungen sollen erst beim Lesen wahrgenommen werden, sich aber ansonsten in den Text eingliedern.
Die typische Form einer integrierten Auszeichung ist die \emph{kursive} Schrift.
Aktive Auszeichnungen sollen dagegen sofort beim Betrachten der Seite auf"|fallen.
Der wichtigste Vertreter ist hier die \textbf{fette} Schrift.
In wissenschaftlichen Arbeiten werden vorwiegend integrierte Auszeichnungen benutzt.

Grundsätzlich sollte bei Auszeichnungen immer nur \emph{ein} Attribut geändert werden, also \emph{nicht} gleichzeitig \emph{\underline{\textbf{fett, kursiv und
unterstrichen}}}.
Programmcode und Befehle setzt man üblicherweise in \texttt{Schreib\-ma\-schin\-en\-schrift}, Namen gelegentlich in \textsc{Kapitälchen}.
Für manche Bezeichnungen kommt eine \textsf{\textbf{fette serifenlose Schrift}} in Frage.
Hier müssen ausnahmsweise \emph{zwei} Attribute geändert werden, da sich die \textsf{serifenlose Schrift} zu wenig vom übrigen Text abhebt.

\glqq Anführungszeichen\grqq\ (siehe \secref{sec:Sonderzeichen}) sind sparsam zu verwenden, \zB bei umgangssprachlichen Begriffen oder wörtlichen Zitaten.
\underline{Unterstreichen} und \mbox{S\,p\,e\,r\,r\,e\,n} sollen überhaupt nicht benutzt werden.
Es ist wichtig, sich zu Beginn der Arbeit zu überlegen, welche Begriffe in welcher Schrift gesetzt werden, und dies konsequent einzuhalten.
So können \bspw die Namen von Autoren, wie in \name{Dirac}'sche Deltafunktion, einheitlich in Kapitälchen gesetzt werden.



\subsubsection*{\LaTeX}

\tabref{tab:Hervorhebungen} fasst die \LaTeX-Befehle zu den wichtigsten Auszeichnungen und Hervorhebungen zusammen.

Eine Besonderheit ist der Befehl \verb!\verb|...|!. Mit diesem wird der zwischen den senkrechten Strichen eingeschlossene Text Zeichen für Zeichen, \dah ohne, dass er von \LaTeX{} als Code interpretiert wird, in der Schreibmaschinenschrift geschrieben.
Man kann auch anstelle der senkrechten Striche ein anderes Zeichen verwenden, wenn der darzustellende Text gerade selber einen senkrechten Strich enthält.

\begin{table}
  \caption{Auszeichnungen und Hervorhebungen.}
  \label{tab:Hervorhebungen}
  \setlength{\tabcolsep}{1em}\centering
  \begin{tabular}{lll}
    \toprule
          &  Beispiel        &  Eingabe\\
    \midrule
    kursive Schrift      &  \emph{Beispiel}      &  \verb*|\emph{Beispiel}| \\
    & & oder \verb*|\textit{Beispiel}|\\
    fette Schrift      &  \textbf{Beispiel}      &  \verb*|\textbf{Beispiel}|\\
    Schreibmaschinenschrift      &  \texttt{Beispiel}      &  \verb*|\texttt{Beispiel}| \\
    & & oder \verb*!\verb|Beispiel|! \\
    serifenlose Schrift  &  \textsf{Beispiel}      &  \verb*|\textsf{Beispiel}|\\
    fette, serifenlose Schrift  &  \textsf{\textbf{Beispiel}}      &  \verb*|\textsf{\textbf{Beispiel}}|\\
    Kapitälchen      &  \textsc{Beispiel}      &  \verb*|\textsc{Beispiel}|\\
    Anführungszeichen (deutsch) &  "`Beispiel"'      &  \verb*|"`Beispiel"'| \\
    & &  oder \verb*|\glqq Beispiel\grqq|\\
    Anführungszeichen (englisch) &  ``Beispiel''    &  \verb*|``Beispiel''|\\
    \bottomrule
  \end{tabular}
\end{table}



\section{Einbinden von Quellcode}
Wird Quellcode (\Matlab{}, C, \ldots) in der Arbeit angegeben, ist grundsätzlich eine Monospace-Schriftart zu verwenden, da nur so die Lesbarkeit des Codes
gewährleistet werden kann.

Manchmal kann es auch erforderlich sein, die Schrift zu verkleinern oder notfalls sogar die Seiten im Querformat zu beschreiben.
Im laufenden Text sollten nur kleinere Code-Fragmente abgedruckt sein, längere Programme gehören, wenn diese überhaupt abgedruckt werden sollen, grundsätzlich in den Anhang.



\subsubsection*{\LaTeX}

Die einfachste Möglichkeit zum Einbinden von Quellcode ist die \verb|verbatim|- oder \verb|verbatim*|-Umgebung.
Der Code wird in Schreibmaschinenschrift \emph{exakt} (inklusive aller Leer- und Sonderzeichen) so wiedergegeben, wie er im \LaTeX-Quelltext steht.
Die \verb|*|-Variante der Umgebung druckt die Leerzeichen als \verb*| |, was teilweise hilfreich ist.

Komfortablere Möglichkeiten bietet das \verb|listings|-Paket, \zB Syntax-Highlighting mit verschiedenen Schriften oder das Einbinden externer Dateien.


\section{Abstände und Sonderzeichen}
\label{sec:Sonderzeichen}

Bei fest verbundenen Begriffen benutzt man ein \emph{festes} Leerzeichens, \zB bei "`Dr.~Müller"' "`3~Uhr"', das weder umgebrochen noch gedehnt werden kann.

Bei zusammengesetzten Abkürzungen, beispielsweise \glqq\dah{}\grqq, \glqq\ua{}\grqq\ oder \glqq\zB{}\grqq, wird ein \emph{kleiner} Zwischenraum verwendet.
Der kleine Zwischenraum steht auch zwischen Zahl und Einheit bei physikalischen Größen, siehe \tabref{tab:Sonderzeichen}.

Unterschiede sind auch bei den \emph{Strichen} zu beachten.
Der \emph{Bindestrich} "`-"' steht bei zusammengesetzten Wörtern oder Trennungen und wird ohne zusätzlichen Zwischenraum benutzt.
Der \emph{Gedankenstrich} "`--"' steht bei eingeschobenen Satzteilen und als "`Bis-Strich"'.
Als Gedankenstrich wird er immer mit einem Leerzeichen davor und dahinter benutzt, als \glqq Bis-Strich\grqq\ ohne Leerzeichen. Im Englischen ist der Gedankenstrich länger, "`---"', und wird ohne Leerzeichen davor und dahinter verwendet.
In mathematischen Ausdrücken wird das \emph{Minuszeichen} "`$-$"' verwendet.

Es ist darauf zu achten, im Deutschen die "`deutschen Anführungszeichen"' und im Englischen die ``englischen Anführungszeichen'' zu verwenden.
Dies richtet sich nach der Sprache, in der die Ausarbeitung abgefasst wird, nicht nach der Sprache der hervorgehobenen Worte.


\subsubsection*{\LaTeX}

\LaTeX\ interpretiert ein oder mehrere aufeinander folgende Leerzeichen im Quelltext als normalen Wortzwischenraum.
Nach Befehlen wird es jedoch ignoriert, da es dort nur das Ende des Befehls kennzeichnet.
Soll \zB in dem Satz \glqq\TeX\ ist toll!\grqq\ nach \glqq\TeX\grqq\ ein Leerzeichen erscheinen, dann muss im Quelltext entweder \verb*|\TeX\ | oder \verb|\TeX{}| geschrieben werden.
Im ersten Fall wird durch \verb*|\ | ein Leerzeichen erzwungen, im zweiten Fall wird die leere Umgebung \verb|{}| benutzt, um den Befehl \verb|\TeX| zu beenden.

Ein \emph{festes} Leerzeichen wird durch eine Tilde erzeugt, \zB \verb|Dr.~Müller| oder \verb|3~Uhr|.
Ein \emph{kleiner} Zwischenraum wird durch \verb|\,| gesetzt, \zB in \verb*|d.\,h.\ | oder \verb*|z.\,B.\ |.
Hinter dem zweiten Punkt sollte wie hier gemacht auch ein \verb*|\ | stehen, damit dieser nicht als Satzende interpretiert wird.
(Das beeinflusst die automatische Anpassung der Wortzwischenräume zum Erreichen des Blocksatzes.)

Im \emph{mathematischen} Modus wird das Komma als Aufzählungszeichen interpretiert und dahinter ein kleiner Abstand eingefügt.
Dies ist jedoch problematisch, da das Komma im Deutschen auch als \emph{Dezimal}komma verwendet wird.
Um den zusätzlichen Abstand zu unterdrücken schreibt man \zB \verb|$2{,}5x$| für \glqq $2{,}5x$\grqq, oder man verwendet das \LaTeX-Paket \verb|icomma|, welches dies -- in der Regel -- automatisch erkennt.

Im normalen Textmodus wird mit \verb|-| ein einfacher Bindestrich erzeugt.
Den längeren Gedankenstrich erhält man mit \verb|--|, den englischen Gedankenstrich mit \verb|---|. Im Mathemodus wird mit \verb|-| automatisch das längere Minuszeichen gesetzt.

Die deutschen Anführungszeichen werden mit \verb|\glqq| und \verb|\grqq{}| \bzw \verb*|\grqq\ | gesetzt.
Wie oben erläutert, muss der Befehl \verb|\grqq| mit \verb*|\ | oder mit \verb|{}| abgeschlossen werden, falls danach ein Leerzeichen folgen soll.
Alternativ können die deutschen Anführungszeichen auch mit \verb|"`| und \verb|"'| gesetzt werden.
Hier ent"|fällt die Notwendigkeit, das schließende Zeichen besonders abzuschließen.%
\footnote{Es ist mit der hier verwendeten Schriftart leider kaum zu erkennen:
  Das öffnende deutsche Anführungszeichen wird über "`Doppeltes Anführungszeichen"' + "`oberes Akzentzeichen (Gravis)"' eingegeben, das schließende über "`doppeltes Anführungszeichen"' + "`einfaches Anführungszeichen"'.
  Die englischen Anführungszeichen über "`oberes Akzentzeichen"' + "`oberes Akzentzeichen"' und "`einfaches Anführungszeichen"' + "`einfaches Anführungszeichen"'.}

Die englischen Anführungszeichen werden mit \verb|``| und \verb|''| gesetzt.

\tabref{tab:Sonderzeichen} zeigt Beispiele zu den hier beschriebenen Abständen und Zeichen.

\begin{table}
  \caption{Die wichtigsten Abstände und Sonderzeichen.}
  \label{tab:Sonderzeichen}
  \setlength{\tabcolsep}{1em}\centering
  \begin{tabular}{lll}
  \toprule
    Bezeichnung      &  Beispiel        &  Eingabe\\
  \midrule
    Leerzeichen      &  \TeX\ ist toll!      &  \verb*|\TeX\ ist toll!|\\
              &              &  \verb*|\TeX{} ist toll!|\\
    festes Leerzeichen  &  Dr.~Müller        &  \verb|Dr.~Müller|\\
    kleines Leerzeichen  &  d.\,h.\ 3,5\,km      &  \verb*|d.\,h.\ 3,5\,km|\\
    Bindestrich      &  \TeX-Datei        &  \verb|\TeX-Datei|\\
    Gedankenstrich    &  S.~153--165        &  \verb|S.~153--165|\\
    Minuszeichen    &  $y=5x-2$        &  \verb|$y=5x-2$|\\
    Anführungszeichen (deutsch)  &  \glqq Beispiel\grqq\  &  \verb*|\glqq Beispiel\grqq\ |\\
            &              &  oder \verb*|\glqq Beispiel\grqq{}| \\
            & & oder \verb*|"`Beispiel"'| \\
  Anführungszeichen (englisch)  &  ``Beispiel''  &  \verb*|``Beispiel''|\\
  \bottomrule
  \end{tabular}
\end{table}



\section{Stilistische Punkte}

Über den Stil lässt sich natürlich immer streiten, spielt dabei ja auch der persönliche Geschmack eine große Rolle.
Jedoch gibt es ein paar Anhaltspunkte für die Abfassung einer wissenschaftlichen Abschlussarbeit, an denen man sich orientieren sollte.

Der wesentliche Punkt ist, dass Klarheit und Verständlichkeit vor sprachlicher Schönheit kommen.
Sie sollten also nicht nach schönen oder gar lyrischen Formulierungen suchen, sondern nach einfachen.
Sie müssen keine Fremdwörter verwenden, wenn es für das, was Sie ausdrücken möchten, auch einfache Wörter in der Sprache existieren, in der Sie die Ausarbeitung abfassen.
Sie sollten darauf achten, dass die Sätze nicht zu lang werden.
Das bedeutet keinesfalls, dass Sie keine Nebensätze verwenden dürfen, aber es sollte auch nicht so sein, dass sich ein einzelner Satz über viele Zeilen oder ganze Absätze zieht.
Auch wenn für Sie solche Sätze (zunächst) einfach zu lesen sind, ist es dies für andere Personen, die Ihre Arbeit lesen, nicht.
Und es sollte auch Ihnen helfen, Ihre Gedanken in kleinere zusammenhängende Abschnitte zu zerlegen.

Es sollte bei einer wissenschaftlichen Abfassung auch nicht das Ziel sein, möglichst variantenreich zu schreiben, indem ständig andere Wörter und Ausdrücke für die gleiche Sache verwendet werden.
Hier ist es günstiger, das Gleiche auch immer gleich zu bezeichnen, auch wenn man das Gefühl hat, dabei "`langweilig"' zu klingen.

Gleichungen und (knappere) Aufzählungen sollten als Bestandteil des Textes aufgefasst und idealerweise in die Sätze eingearbeitet werden.
Zum einen wird dies meist als "`eleganter"' empfunden, zum anderen macht es auch das Lesen des Textes einfacher.
Schauen Sie hierzu auch einmal in Ihre Lehrbücher.


\chapter{Beurteilung}
\label{cha:Beurteilung}

Die Beurteilung der Arbeit erfolgt nach einem Schema, welches die Hauptkriterien
\begin{itemize}
    \item Arbeitsstil,
    \item Ergebnisse,
    \item Ausarbeitung und
    \item Abschlussvotrag
\end{itemize}
umfasst.
Im Folgenden sind kurz die wesentlichen Aspekte, die darunter jeweils betrachtet werden, beschrieben.


\paragraph{Arbeitsstil}
Hierunter fällt die Bewertung des (hoffentlich) systematischen und methodischen Vorgehens.
Dies umfasst neben der sinnvollen zeitlichen Vorgehensweise auch die Literaturrecherche zu Beginn der Arbeit sowie die Kooperation und Kommunikation mit dem Betreuer.
Daneben spielt auch die Selbständigkeit sowie das fachliche Verständnis eine große Rolle.

Bei Gruppenarbeiten wird auch die Zusammenarbeit der Gruppe bewertet.
Dies meint nicht unbedingt, wie harmonisch eine Gruppe zusammenarbeitet, sondern ob die Aufgabe in sinnvolle Pakete zerlegt wurde, so dass durch die Bearbeitung in der Gruppe ein Mehrwert im Vergleich zu Einzelarbeiten erreicht wurde.


\paragraph{Ergebnisse}
Hierbei wird bewertet, inwieweit die Arbeit vollständig ist und welche Qualität die Ergebnisse haben.
\Dah sind alle geforderten Punkte (und \ggf naheliegende Aspekte) berücksichtigt und in einer solchen Form bearbeitet, dass diese direkt weiterverwendet werden können?


\paragraph{Ausarbeitung}
Ein wesentliches formales Kriterium an dieser Stelle ist die korrekte und vollständige Zitierweise.
Schwerwiegende Fehler an dieser Stelle (Plagiate) führen zu einem Nicht-Bestehen der gesamten Arbeit!

Darüber hinaus sind die wesentlichen Kriterien die Gliederung der Arbeit sowie die Vollständigkeit und Verständlichkeit.
\Dah ist der Aufbau und Argumentationsweg der Arbeit nachvollziehbar, sind die Begründungen schlüssig?

Die sprachliche Ausdrucksweise (über die Verständlichkeit hinaus) und Korrektheit, sowie die prinzipielle Sorgfalt der Gestaltung (im Wesentlichen Plots und Zeichnungen) werden hierbei auch berücksichtigt, aber mit einem geringeren Gewicht.


\paragraph{Abschlussvortag}
Auch hier wird die Gliederung der Folien bewertet.
Wurde sich auf das wichtigste beschränkt und folgt der Vortrag einem roten Faden?
Auch die Gestaltung der Folien spielt hier eine Rolle bei der Bewertung.
Der Vortragsstil wird ebenfalls beachtet, jedoch mit einem etwas geringeren Gewicht.

Neben dem eigentlichen Vortrag spielt hier auch die Diskussion eine wesentliche Rolle bei der Bewertung.
\Dah wurden die Fragen korrekt erfasst und entsprechend eine überzeugende Antwort darauf gegeben?



\paragraph{Vorkorrektur und Probevortrag}
Wir bieten jeweils an, die schriftliche Ausarbeitung vor der Abgabe einmal zu lesen und Korrekturen \bzw Anmerkungen dazu zu geben.
Ebenso bieten wir an, einen Probevortrag vor dem eigentlichen Seminartermin zu halten.
Hierzu ist klarzustellen, dass sich die Bewertung der Ausarbeitung sowie des Aufbau des Vortrages sich schwerpunktmäßig auf diese Vorabversion \bzw den Probevortrag bezieht, da wir letztlich nicht unsere eigenen Korrekturen bewerten wollen.
Natürlich wird aber eine entsprechende Umsetzung etwaiger Hinweise bei der Bewertung berücksichtigt, so dass eine Vorabgabe und ein Probevortrag immer vorteilhaft ist.
(\Dah eine endgültige Abgabe der Arbeit mit größeren Mängeln ist immer "`schlimmer"', als wenn diese noch nach der Vorkorrektur verbessert werden.)



% =================================================================================
% Anhang
% =================================================================================
\appendix % Damit wird der Anhang begonnen. Die Kapitel werden ab jetzt mit Buchstaben nummeriert

\chapter{Checkliste}
\label{cha:Checkliste}

Aktuell ist von der Arbeit eine elektronische Version der Arbeit im PDF/A Format in \href{https://tubama.ulb.tu-darmstadt.de/}{\textsc{TUbama}} abzugeben.
Zusätzlich ist die Arbeit in elektronischer Form (PDF) direkt an den Betreuer zu schicken.

In Absprache mit dem Betreuer kann auch ein gedrucktes Exemplar (doppelseitig; gebunden mit dunkelblauem oder schwarzem Karton; vorne Klarsichtfolie; Klebe- oder Leimbindung, keine Spiralbindung, kein Hardcover) eingefordert werden.

Zusätzlich ist mit dem Betreuer abzusprechen, welche weitere Daten in elektronischer Form abzugeben sind (Quelltext wie Matlab-Dateien und Simulinkmodelle, technische Zeichnungen und CAD-Modelle, Literatur).

Zur Abgabe der Arbeit sollten folgende Punkte überprüft werden:

% Der folgende Befehl ändert die Aufzählungszeichen in Kästen. Diese Definition
% ist im gesamten restlichen Dokument wirksam.

\renewcommand{\labelitemi}{%
  \setlength{\fboxsep}{.6ex}\raisebox{.7ex}{\framebox{}}}

\textbf{Im Quelltext}
\begin{itemize}
  \item Wurden alle der Literatur entnommenen Stellen mit Literaturverweisen belegt?
  \item keine wörtlichen Zitate, falls doch müssen diese in Anführungszeichen!
  \item Alle Literaturangaben im Literaturverzeichnis vollständig? (Autor(en), Titel, Verlag/Journal/Konferenzband/\etc, Jahr, \usw)
  \item Alle Seitenangaben im Inhaltsverzeichnis korrekt?
  \item Keine einzelnen Abschnitte/Unterabschnitte im Inhaltsverzeichnis?
  \item Prägnante Kapitelnamen? (keine ganzen Sätze!)
  \item Alle Bilder haben Bild\emph{unterschriften}, alle Tabellen haben \emph{Überschriften}?
  \item Wurden die Hinweise aus \charef{cha:Hinweise_Allgemein} berücksichtigt (allgemeiner Aufbau, Abstände, Sonderzeichen, \ldots)?
  \item Sind die Bilder gut erkennbar und alle Elemente beschriftet?
    Passt die Schriftgröße in den Bildern zum Text?
    Sitzen die Gleitobjekte an der richtigen Stelle?
  \item Treten beim Aufruf von \TeX\ bzw.\ pdf\TeX\ Fehler oder Warnungen auf?
    Sind \emph{alle} Bilder, Tabellen und Literaturstellen im  Text zitiert?
    Sind falsche oder doppelte Referenzen vorhanden?
    Dies lässt sich anhand der Log-Datei feststellen.
\end{itemize}

\textbf{In der PDF-Datei}
\begin{itemize}
  \item Die Datei ist im \textbf{doppelseitigen} Layout mit Hypertext-Elementen zu erstellen; die Optionen \verb|draft|, \verb|oneside| und \verb|nohyperref| sind \textbf{nicht} aktiviert.
    Die Seitengröße des PDF-Dokuments überprüfen:
    210\,mm $\times$ 297\,mm (DIN A4).
  \item Sind die PDF-Bookmarks und Seitenzahlen vorhanden?
    Zumindest im Vorspann sollte überprüft werden, ob die Bookmarks auf die richtige Seite verweisen.
  \item Sind die PDF-Infofelder (in Acrobat Reader: Datei $\rightarrow$ Dokumenteigenschaften) richtig eingetragen?
    Notfalls mit \verb|\hypersetup{...}| korrigieren.
  \item In den Bookmarks und Infofeldern können nicht alle Zeichen dargestellt werden.
    In einem solchen Fall \zB\ \verb|\texorpdfstring| (aus \verb|hyperref|) verwenden.
  \item Ist die für \textsc{TUbama} erzeugte Datei eine gültige PDF/A Datei?
\end{itemize}

\textbf{Im Ausdruck (wenn gefordert)}
\begin{itemize}
  \item Die Arbeit ist \textbf{doppelseitig}, vorzugsweise schwarzweiß auf einem Laserdrucker auszudrucken.
    Sind alle Grafiken gut zu erkennen (Farbe, Linienstärke, \etc)?
    Wurden alle Sonderzeichen korrekt gedruckt?
  \item Wurde eine Klebe- oder Leimbindung zum Binden verwendet?
  \item Beim Ausdrucken aus dem Acrobat Reader darf die Seitenanpassung \textbf{nicht} aktiviert sein, da der Textblock sonst verkleinert wird.
    Ränder müssen zum Binden geeignet eingestellt sein.
  %\item Selbstständigkeitserklärung unterschreiben!
  %\item Es ist eine zusätzliche eigenständige Selbstständigkeitserklärung auszudrucken und zu unterschreiben.
\end{itemize}



\chapter{Programme zur Erstellung von Grafiken}
\label{cha:Anhang-Grafiken}

\section{Vektorgrafiken}

Vektorgrafiken bestehen aus geometrischen Formen, deren Beschreibung unabhängig von der Auflösung ist. Sie sind \zB\ für Diagramme und mathematische Plots sinnvoll und lassen sich aus vielen Programmen direkt als EPS oder PDF speichern.
Es ist besonders darauf zu achten, dass die Schriften und Strichstärken zum Rest des Dokuments passen.

\renewcommand{\labelitemi}{$\bullet$}

\begin{itemize}
  \item \LaTeX\ selbst bietet mit den Paketen \texttt{TikZ} und \texttt{pgfplots} zwei sehr mächtige Werkzeuge, um direkt in \LaTeX\ Vektorgrafiken, wie zum Beispiel Blockschaltbilder oder Plots zu erzeugen.
    Schriftart und -größe sind automatisch identisch mit dem übrigen laufenden Text, so dass hier keine Anpassungen mehr vorgenommen werden müssen.
    \begin{itemize}
      \item \texttt{TikZ} eignet sich hervorragend für Blockschaltbilder.
      \item Mit \texttt{pgfplots} können aus einer \texttt{.txt}-Datei mit den Variablenwerten Plots direkt in \LaTeX\ erzeugt werden, so dass sich ein einheitliches Gesamtbild ergibt.
        Die benötigte \texttt{.txt}-Datei lässt sich mit \Matlab\ einfach erzeugen.
        Achsenbeschriftungen, Legenden, \etc können mit \Matlab-ähnlichen Befehlen einfach hinzugefügt werden.
      \item Für \Matlab\ gibt es die über den Mathworks File Exchange zu ladende Funktion \verb|matlab2tikz|, mit der Matlab-Plots in das \texttt{pgfplots}-Format konvertiert werden können.
        Die dazugehörende Funktion \verb|cleanfigure| sorgt dabei für eine effektive Reduzierung der Plotpunkte, was die Bildgröße (Speicherbedarf) deutlich reduzieren kann.
    \end{itemize}
    Die Dokumentationen zu \texttt{TikZ}~\cite{TikZ} und \texttt{pgfplots}~\cite{pgfplots} sind sehr ausführlich und mit vielen Beispielen leicht verständlich erklärt.
  \item Alternativ können Sie auch andere Zeichenprogramme verwenden, bei denen Sie "`wie gewohnt"' mit der Maus arbeiten können.
    Ein solches Zeichenprogramm ist \zB Inkscape.
    Inkscape ermöglicht auch den Export der Zeichungen in ein \LaTeX-Format.
  \item Wenn Sie andere Programme verwenden, müssen Sie darauf achten, dass das Ergebnis möglichst als PDF exportiert werden kann, um es als Vektorgrafik einbinden zu können.
    Ggf.\ ist der Umweg über den Export aus dem verwendeten Programm im EPS-Format und der anschließenden Konvertierung der EPS-Datei in ein PDF möglich.
    Hierfür eignen sich der Acrobat Distiller oder Ghostscript in Verbindung mit EPSTOPDF (das in allen gängigen TeX-Distributionen enthalten ist).
\end{itemize}



\section{Pixelgrafiken}

Pixelgrafiken besitzen eine feste Anzahl von Bildpunkten, \dah ihre Auflösung hängt von der Größe der Darstellung ab.
Man benutzt Pixelformate \zB für Fotos, Screenshots oder eingescannte Grafiken.
Hierbei ist die Wahl der Auflösung besonders wichtig.
Die Bilder sollen einerseits eine gute Druckqualität ergeben, andererseits aber auch eine zügige Bildschirmdarstellung und kleine Dateigröße ermöglichen.
Beim Scannen ist außerdem zu beachten, dass sich die Auflösung ändert, wenn die Grafiken nicht in Originalgröße eingebunden werden.
\begin{itemize}
\item Fotos liegen in der Regel im JPEG-Format vor; eine Auflösung von $100$--$\valunit{150}{dpi}$ ist häufig bereits ausreichend.
  pdf\LaTeX\ und LuaLaTeX können JPEG-Grafiken direkt einlesen.
\item Sonstige Farb- oder Graustufen-Grafiken (insb.\ wenn sie \glqq harte\grqq\ Farbübergänge besitzen) werden am besten im PNG-Format gespeichert.
  Dieses Format wendet keine verlustbehaftete Komprimierung an.
  Damit wird bei Screenshots oder bei anderen Grafiken mit harten Kanten verhindert, dass diese Kanten "`verschwimmen"'.
  Die richtige Wahl der Farbtiefe hat großen Einfluss auf die spätere Dateigröße.
  Bzgl.\ der Auflösung gibt es hier keine allgemeine Regel; oft liegt sie bereits fest (\zB bei Screenshots).
  pdf\LaTeX\ und LuaLaTeX können PNG-Grafiken direkt verarbeiten.
\item Schwarzweiße Strichzeichnungen müssen in relativ hoher Auflösung ($\ge \valunit{300}{dpi}$) vorliegen, um beim Drucken eine ausreichende Qualität zu gewährleisten.
  Diese sollten im png-Format gespeichert werden.
\end{itemize}


% =================================================================================
% Literaturverzeichnis
% =================================================================================
\cleardoublepage        % Auf eine leere Seite einfügen
\phantomsection         % Für Aufnahme ins Inhaltsverzeichnis
\addcontentsline{toc}{chapter}{\bibname}  % In Inhaltsverzeichnis von
                                          % Dokument und pdf aufnehmen
\printbibliography
% =================================================================================


\end{document}
